\documentclass[11pt]{article}
\usepackage[T1]{fontenc}
\usepackage[utf8]{inputenc}
\usepackage{lmodern,amsmath,amssymb,amsthm,mathtools}
\usepackage[margin=1in]{geometry}
\usepackage{xurl}
\usepackage[hidelinks]{hyperref}
\usepackage{microtype}
\hypersetup{pdftitle={Band edges of the almost Mathieu operator at Liouville scales},pdfauthor={Seth Douglas and Nidhal Mghirbi}}
\newcommand{\C}{\mathbb C}
\newcommand{\N}{\mathbb N}
\newcommand{\Z}{\mathbb Z}
\newcommand{\Q}{\mathbb Q}
\newcommand{\R}{\mathbb R}
\newcommand{\T}{\mathbb T}
\newcommand{\un}{\mathbf 1}
\newcommand{\tr}{\operatorname{tr}}
\newcommand{\ad}{\operatorname{ad}}
\newcommand{\dist}{\operatorname{dist}}
\newcommand{\supp}{\operatorname{supp}}
\newcommand{\Lip}{\operatorname{Lip}}
\newcommand{\Ga}{G_\alpha}
\newcommand{\au}{\alpha_u^+}
\newcommand{\aunif}{\alpha_u^{\mathrm{unif}}}
\newcommand{\eps}{\varepsilon}
\newcommand{\om}{\omega}
\newcommand{\omt}{\tilde\omega}
\newcommand{\omh}{\hat\omega}
\newcommand{\FU}{F_{\mathrm U}}
\newtheorem{theorem}{Theorem}[section]
\newtheorem{lemma}[theorem]{Lemma}
\newtheorem{proposition}[theorem]{Proposition}
\newtheorem{corollary}[theorem]{Corollary}
\newtheorem{conjecture}[theorem]{Conjecture}
\newtheorem{maintheorem}{Theorem}
\renewcommand{\themaintheorem}{\Alph{maintheorem}}
\newtheorem{maincorollary}[maintheorem]{Corollary}
\theoremstyle{definition}
\newtheorem{definition}[theorem]{Definition}
\newtheorem{remark}[theorem]{Remark}
\newtheorem{example}[theorem]{Example}
\theoremstyle{plain}
\newtheorem*{classicalinput}{Classical input}
\newtheorem*{classicalidentity}{Classical identity}
\newtheorem*{classicalproposition}{Proposition}
\usepackage{framed}
\newenvironment{inputbox}{\begin{framed}\vspace{-0.7\baselineskip}}{\vspace{-0.7\baselineskip}\end{framed}}
\setlength{\emergencystretch}{2em}
\raggedbottom

\usepackage{graphicx,booktabs,longtable,array}
\newcommand{\Sig}{\Sigma_{\lambda,\alpha}}
\newcommand{\Ran}{\operatorname{Ran}}
\newcommand{\sgn}{\operatorname{sgn}}
\newcommand{\Op}{\operatorname{Op}}
\newcommand{\cF}{\mathcal F}
\newcommand{\cV}{\mathcal V}
\newcommand{\cU}{\mathcal U}
\newcommand{\cI}{\mathcal I}
\newcommand{\hM}{\mathsf h}
\newcommand{\EE}{\mathcal E_q}
\newcommand{\neig}{\mathcal N}
\newcommand{\inner}[2]{\langle #1,#2\rangle}
\newcommand{\norm}[1]{\lVert #1\rVert}
\newcommand{\abs}[1]{\lvert #1\rvert}
\title{Band edges of the almost Mathieu operator\\ at Liouville scales}
\author{Seth Douglas \and Nidhal Mghirbi}
\date{Version 1.0.0, November 2026}
\begin{document}
\maketitle
\begin{center}\small Both authors contributed equally.\end{center}
\begin{abstract}
Can a supercritical almost Mathieu operator have gaps that are only polynomially small in their integer label? We prove that it does whenever its frequency has finite Liouville exponent $\beta(\alpha)>\tfrac12\log\lambda$. The proof uses the top edge of a rational approximant and imposes no parity condition on its denominator. At a convergent $p_n/q_n$, write $\delta_n=q_n\alpha-p_n$, $b_n=-q_n^{-1}\log|\delta_n|$ and $L=\log\lambda$. We construct gaps with labels $m=-Dq_n\operatorname{sgn}\delta_n$ and width at least $\exp(-(b_n+L/2+o(1))q_n)$, uniformly over an explicit growing range of $D$. This gives the power $(\beta'+L/2)/(\beta'-L/2)$ below $\beta=3L/2$, for every $L/2<\beta'<\beta$, and $(\beta+L/2)/L$ at and above that value, with arbitrary positive slack. Hence the exponential gap-decay rate vanishes. A separate surplus argument includes $\beta=\infty$. The reduction combines periodic Green-function decay, translated Wannier frames, self-adjoint comparison operators and exact finite-box counts. The orbital displacement is the classical Rammal--Wilkinson correction. The quantitative contribution is their uniform combination at one Liouville scale with growing labels. We also give the rational surplus argument, a fixed-curvature Mathieu limit, and a phase diagram that separates established results from conjectures. No computation enters a proof.
\end{abstract}
\noindent\textbf{AI disclosure.} AI models assisted the exploration, checking and drafting. The authors are responsible for the claims; the process is described in Section~\ref{sec:discussion}.

{\small\tableofcontents}
\section{Gap decay on both sides of the Liouville threshold}\label{sec:intro}
\paragraph{NEW.}
For a supercritical almost Mathieu operator, positive Liouville exponent changes the gap-decay picture even below the threshold $2\beta=\log\lambda$. The same line is the spectral transition at the completely resonant phases, $2\theta\in\alpha\Z+\Z$: Liu proved Anderson localization there for $\log\lambda>2\beta$~\cite{Liu2022} and, by a doubled Gordon threshold, purely singular continuous spectrum for $1<\lambda<e^{2\beta}$~\cite{Liu2026}. We prove an upper bound on the normalized decay rate, $\rho\le1-2\beta/\log\lambda$, below that threshold. Together with the polynomial lower subsequences established in version~1.0.0 of this paper~\cite{PaperDv1}, this gives quantitative gap-width results on both sides. The equality line and the matching lower bound below it remain open.

\paragraph{KNOWN (normalization and gap labels).}
We use hopping one:
\begin{equation}\label{eq:H}
(H_{\lambda,\alpha,\theta}u)(j)=u(j+1)+u(j-1)
 +2\lambda\cos(2\pi(j\alpha+\theta))u(j).
\end{equation}
Throughout, $\lambda>1$, $L=\log\lambda$, and $\alpha$ is irrational. Its convergents are $p_n/q_n$, and
\begin{equation}\label{eq:arithmetic}
\delta_n=q_n\alpha-p_n,\quad b_n=-\frac{\log|\delta_n|}{q_n},
\quad \beta=\limsup_n\frac{\log q_{n+1}}{q_n}=\limsup_n b_n.
\end{equation}
The common spectrum is $\Sig$. The integrated density of states $N$ labels a bounded gap $G_m$ by $N|_{G_m}=\{m\alpha\}$, $m\ne0$; see~\cite{AvronSimon,AYZdry}. We normalize the gap exponent as
\begin{equation}\label{eq:rho}
\rho(\lambda,\alpha)=\liminf_{|m|\to\infty}
 \frac{-\log|G_m|}{L|m|},\qquad x=\frac{2\beta}{L}.
\end{equation}
A closed gap contributes $+\infty$. Bounded spectral diameter gives $\rho\ge0$.

\paragraph{KNOWN (rational scales).}
Let $P_n$ be the monic Chambers polynomial and $E_n$ the top of the rational phase union. Set
\begin{equation}\label{eq:scales}
\Lambda_n=\lambda^{q_n},\qquad
\hbar_n=2\pi q_n\Lambda_n^{1/2}|\delta_n|,\qquad
\kappa_n=\hbar_n^2.
\end{equation}
The derivative $P_n'(E_n)$ is positive; see Lemma~\ref{lem:convex}.
\begin{figure}[ht]
\centering\small
\fbox{\begin{minipage}{.93\linewidth}
\centering
Rational top band\quad$\longrightarrow$\quad Translated wells\quad$\longrightarrow$\quad Common ladder\\[4pt]
$P(E_*)=2\Lambda+2$\qquad
one well per $|\delta|^{-1}$ sites\qquad
spacing of order $\hbar/P_*'$\\[7pt]
$D$ levels above $E$ in every well\quad$\Longrightarrow$\quad
$N(E)=1-D|\delta|$\\[3pt]
$\Longrightarrow$\quad a gap with label $m=-Dq\sgn\delta$.
\end{minipage}}
\caption{The single-scale mechanism, with the scale index omitted and $P_*'=P'(E_*)$. The common well count is proved before the infinite-volume IDS limit.}\label{fig:mechanism}
\end{figure}


\paragraph{NEW.}
The following is the new result of version~2.0.0.
\begin{maintheorem}[Fixed-label-multiple gaps below the threshold]\label{thm:below}
Suppose $0<\beta<L/2$. Fix an integer $D\ge1$. There is $c_D>0$, depending only on $D$, such that, on any subsequence of convergents with $b_n\to\beta$, for all sufficiently large $n$,
\begin{equation}\label{eq:belowwidth}
 |G_{-Dq_n\sgn\delta_n}|
 \ge \frac{c_D\kappa_n^{1-D}}{P_n'(E_n)}.
\end{equation}
The threshold in $n$ may depend on $\lambda,\alpha,\beta,D$. In particular, for every $\epsilon>0$, eventually on that subsequence,
\begin{equation}\label{eq:belowexponential}
 |G_{-Dq_n\sgn\delta_n}|
 \ge \exp\!\left(-\bigl[(D-1)(L-2b_n)+L+\epsilon\bigr]q_n\right).
\end{equation}
Consequently $\rho(\lambda,\alpha)\le1-x$ for $0<x<1$.
\end{maintheorem}
No uniformity in growing $D$ is asserted here. The order of limits is first $n\to\infty$ at fixed $D$, then $D\to\infty$.

\paragraph{NEW (combination with a KNOWN estimate).}
\begin{maincorollary}[A positive rate strictly below the Diophantine rate]\label{cor:bracket}
There is an absolute constant $C\ge4$ such that for $\lambda>1$ and irrational $\alpha$ with $0<x\le2/C$,
\[
 \frac1C\le\rho(\lambda,\alpha)\le1-x<1.
\]
Here $C$ can be any sufficiently enlarged constant in Liu--Shi's Theorem~1.3~\cite{LiuShi2019}. Its numerical value is not specified.
\end{maincorollary}

\paragraph{KNOWN (version~1.0.0 of this paper).}
We retain the results on the other side, including their proofs.
\begin{maintheorem}[Polynomial lower subsequences]\label{thm:polynomial}
Suppose $L/2<\beta<\infty$.
\begin{enumerate}
\item If $\beta<3L/2$, for every $\beta'\in(L/2,\beta)$ and $\epsilon>0$ there are infinitely many labels $m$ with
\[
 |G_m|\ge |m|^{-\frac{\beta'+L/2}{\beta'-L/2}-\epsilon}.
\]
\item If $3L/2\le\beta<\infty$, for every $\epsilon>0$ there are infinitely many labels $m$ with
\[
 |G_m|\ge |m|^{-\frac{\beta+L/2}{L}-\epsilon}.
\]
\end{enumerate}
In both cases $\rho=0$. The labels can be chosen as $m=-Dq_n\sgn\delta_n$ at the top of the spectrum.
\end{maintheorem}

\paragraph{KNOWN (version~1.0.0 of this paper).}
The scale estimate behind the preceding theorem permits growing $D$.
\begin{maintheorem}[Gaps at one Liouville scale]\label{thm:scale}
There are absolute constants $c_0,c_1>0$ with the following property. Fix $\xi\in(0,L]$. For all sufficiently large $n$ with $b_n\ge L/2+\xi$, every integer
\begin{equation}\label{eq:Dcap}
1\le D\le
\left\lfloor\min\left\{\frac{c_1}{\hbar_n},
                 \Lambda_ne^{-\xi q_n}\right\}\right\rfloor
\end{equation}
has a gap with
\begin{equation}\label{eq:width}
 N|_{G_m}=1-D|\delta_n|,\quad m=-Dq_n\sgn\delta_n,\qquad
 |G_m|\ge\frac{c_0\hbar_n}{P_n'(E_n)}.
\end{equation}
For each $\eta>0$ the width bound is at least
$C_\eta\exp(-(b_n+L/2+\eta)q_n)$ eventually, with $C_\eta>0$ independent of $b_n,D$.
Equivalently it is $\exp(-(b_n+L/2+o(1))q_n)$, uniformly in $b_n$.
The first cap is smaller when $b_n\le3L/2-\xi$ and $n$ is large.
\end{maintheorem}
The change at $3L/2$ records comparison precision, not a claimed spectral transition. The separate surplus argument in Section~\ref{sec:surplus} covers $\beta=\infty$.

\subsection{Relation to earlier work}
\paragraph{KNOWN.}
All noncritical allowed gaps are open by Avila--You--Zhou~\cite[Theorem~1.1]{AYZdry}. Our results concern widths. Liu--Shi~\cite[Theorem~1.3]{LiuShi2019} prove exponential gap upper bounds at sufficiently small dual coupling. Section~\ref{sec:below-literature} proves the label-preserving duality passage and quotes their Remark~1.4(4) exactly. Theorem~\ref{thm:below} is contrary to the expectation in that remark on the nonempty overlap $0<x\le2/C$. It does not contradict their theorem, and it does not prove $\rho=1-x$.

\paragraph{KNOWN IN A SPECIAL CASE (Diophantine frequencies).}
Ge--You--Zhou~\cite[Theorem~1.1]{GYZ} prove $\rho=1$ for Diophantine $\alpha$; earlier estimates are due to Leguil--You--Zhao--Zhou~\cite{LYZZ2017}. These statements do not settle every frequency with $\beta=0$. Avila--You--Zhou~\cite[Theorem~6.2]{AYZdry} also give rational-gap estimates in a different label range. At critical coupling, Krasovsky~\cite[Theorems~3--4]{Krasovsky} obtains polynomial inherited-gap bounds with approximation and parity hypotheses. Polynomial Liouville-gap phenomena themselves are therefore known.

\paragraph{KNOWN (the reduction mechanisms).}
Wilkinson~\cite{Wilkinson1987} develops projected Hamiltonians near rational magnetic flux. Helffer--Sj\"ostrand~\cite[Section~5.1]{HSII} give a rigorous one-step reduction with translated reference wells; their iterated results have additional arithmetic hypotheses. The orbital displacement is the Rammal--Wilkinson correction~\cite{RammalBellissard,BellissardBarelli}. Gram orthogonalization and interaction matrices are standard multiple-well tools~\cite[Sections~2.1--2.2 and 3.2]{HSwells}. The IDS limit is classical~\cite{AvronSimon}.

\paragraph{NEW (the quantitative combination).}
Below the threshold we retain the full energy-dependent self-energy. A specified continuation removes remote physical wells and produces one covariant family. A relative walk estimate resolves gaps of order $\kappa^{1-D}$ at each fixed $D$. A common spectral margin and total Schur error then give an exact integer count in every long box. Above the threshold the retained proof uses common Mathieu windows with growing $D$. Neither argument claims priority for its individual mechanisms.

\subsection{Reading the proof}
\paragraph{NEW / KNOWN (proof organization).}
Section~\ref{sec:below} gives the fixed-$D$ reduction and the Liu--Shi comparison. The earlier above-threshold conclusions and surplus proof remain in Sections~\ref{sec:consequences}--\ref{sec:surplus}. The phase diagram separates the new gap bound from known spectral and transport statements. Appendices~\ref{app:complex}--\ref{app:limits} retain the earlier machinery. Appendices~\ref{app:below-resolvent}--\ref{app:below-assembly} prove the new continuation, self-energy estimates, covariant model theorem and IDS assembly. Appendix~\ref{app:companions} states every companion import with its source.
\label{end:intro}

\section{From a single scale to polynomial gaps}\label{sec:consequences}
\subsection{The arithmetic conversion}
The continued-fraction inequalities give
\begin{equation}\label{eq:CF}
 \frac1{q_{n+1}+q_n}<|\delta_n|<\frac1{q_{n+1}},\qquad
 0<b_n-\frac{\log q_{n+1}}{q_n}<\frac{\log2}{q_n}.
\end{equation}
In particular $\limsup b_n=\beta$. Put
\[
 D_n=\left\lfloor\min\{c_1/\hbar_n,
                        e^{(L-\xi)q_n}\}\right\rfloor,
 \qquad m_n=-D_nq_n\sgn\delta_n.
\]
For $0<\xi<L$, when $b_n\ge L/2+\xi$, both quantities inside the minimum tend to infinity. Rounding therefore changes their logarithm by $o(1)$. For bounded $b_n$,
\begin{align}
 \log|m_n|&=\min\{b_n-L/2,L-\xi\}\,q_n+O(\log q_n),\label{eq:labelsize}\\
 -\log|G_{m_n}|&\le(b_n+L/2+\eta)q_n+O_\eta(1).
 \label{eq:gapsize}
\end{align}
The minimum must be retained near the cap switch. The factor $2\pi q_n$ in $\hbar_n$ shifts its exact location by $O(\log q_n/q_n)$.

\begin{proof}[Proof of Theorem~\ref{thm:polynomial} from Theorem~\ref{thm:scale}]
Suppose first $L/2<\beta<3L/2$, and fix $\beta'\in(L/2,\beta)$. Choose $\xi>0$ smaller than $\beta'-L/2$ and with $\beta+2\xi<3L/2$. Infinitely many $n$ satisfy $b_n\ge\beta'$, and all sufficiently large $n$ satisfy $b_n\le\beta+\xi$. On this subsequence the first cap is active. Dividing~\eqref{eq:gapsize} by~\eqref{eq:labelsize} gives
\[
 \frac{-\log|G_{m_n}|}{\log|m_n|}
 \le\frac{b_n+L/2+\eta}{b_n-L/2}+o(1)
 \le\frac{\beta'+L/2}{\beta'-L/2}
       +\frac{\eta}{\beta'-L/2}+o(1).
\]
Choose $\eta$ sufficiently small and then $n$ sufficiently large to obtain the first assertion.

Now let $3L/2\le\beta<\infty$. Choose $\zeta>0$ small and restrict to infinitely many scales with $b_n>\beta-\zeta$. Eventually $b_n<\beta+\zeta$. Equations~\eqref{eq:labelsize}--\eqref{eq:gapsize} bound the quotient by
\[
 \frac{\beta+\zeta+L/2+\eta}
 {\min\{\beta-\zeta-L/2,L-\xi\}}+o(1).
\]
Let $\zeta,\xi,\eta$ be sufficiently small. This gives any exponent strictly greater than $(\beta+L/2)/L$, including $\beta=3L/2$. Equivalently, diagonalize these fixed-slack choices along the subsequence.

In either case $|m_n|\to\infty$, so there are infinitely many distinct absolute labels. The lower subsequence implies $\rho\le0$. All bounded gaps have length at most $4+4\lambda$, which implies $\rho\ge0$. Hence $\rho=0$.
\end{proof}

\subsection{The count that creates the gap}
Drop the index $n$ for a fixed scale. Let $B_t=H_{\lambda,p/q,t}$, let $E_*$ be the top of its phase union, and write $P_*'=P'(E_*)$. Near phase zero choose a real Wannier function $w_t$ for the top band. Define the even dipole and its energy shift by
\begin{equation}\label{eq:dipole}
 d(t)=\inner{w_t}{XD_tw_t},\qquad
 D_t=\operatorname{diag}v'(mp/q+t),\qquad s(t)=\eps d(t),
 \quad \eps=\alpha-p/q.
\end{equation}
Every well is translated into this same chart. Their common energy origin is $E_*+s(0)$.

Let $\mu_0(\kappa)<\mu_1(\kappa)<\cdots$ be the common Mathieu reference levels constructed in Appendix~\ref{app:mathieu}. Their separation is at least $c_M\hbar$ below a fixed energy less than the separatrix. For
\begin{equation}\label{eq:WD}
 W_D=[\mu_{D-1}+c_M\hbar/4,\ \mu_D-c_M\hbar/4],
 \qquad e=P_*'(E_*+s(0)-E),
\end{equation}
the global reduction proves the following statement. If $D$ satisfies~\eqref{eq:Dcap}, every sufficiently long Dirichlet box whose ends lie midway between wells has exactly $KD$ eigenvalues above $E$ for every $e\in W_D$, where $K$ is its number of wells. The box length is $K/|\delta|+O(1)$.

Taking the box limit gives
\begin{equation}\label{eq:IDSidentity}
 N(E)=1-D|\delta|\qquad(e\in W_D).
\end{equation}
The DOS measure has support $\Sig$. Thus the open interior of the interval in~\eqref{eq:WD}, transported back to energy $E$, contains no spectrum. Its length is at least $c_M\hbar/(2P_*')$. Since $0<D|\delta|<1$, there is spectrum on both sides. Finally,
\[
 (-Dq\sgn\delta)\alpha
 =-D\sgn\delta\,p-D|\delta|,
\]
so the gap label is $-Dq\sgn\delta$. This proves~\eqref{eq:width} once the exact box count is established.

The count is squeezed between two self-adjoint comparison counts, each equal to $KD$. Frame truncation and coupling errors have already been paid inside those operators. There is no residual spectral error to remove by volume normalization. Appendices~\ref{app:global} and~\ref{app:comparison} prove these assertions, completing Theorem~\ref{thm:scale}.

\label{end:consequences}

\section{The rational surplus and the earlier thresholds}\label{sec:surplus}
This argument is independent of the top-band reduction. It also handles $\beta=\infty$. Write $P=P_{p/q}$, $R=2\lambda^q+2$, and
\[
 \cU(p/q)=\{E:|P(E)|\le R\},\qquad
 \cI'(p/q)=\{E:|P(E)|\le R-4\}.
\]
The first set is the rational phase union. The second is a rescaled \emph{dual} phase intersection. It is not the intersection of the supercritical fixed-phase spectra.

\begin{inputbox}
\begin{classicalidentity}[Avron--van Mouche--Simon and Chambers duality]
For every reduced $p/q$,
\begin{align}
 s_q&:=|\cU(p/q)|-4(\lambda-1)\notag\\
 &=|\{E:R-4<|P(E)|\le R\}|\notag\\
 &=\sum_{\text{$2q$ band ends}}
       \int_{R-4}^{R}\frac{dt}{|P'(E(t))|}.\label{eq:surplus}
\end{align}
\end{classicalidentity}
\end{inputbox}
\begin{proof}
Chambers duality is $P_\lambda(E)=\lambda^qP_{1/\lambda}(E/\lambda)$. Avron--van Mouche--Simon~\cite[Theorem~1]{AvMS} give the dual phase-intersection measure $4(1-1/\lambda)$ in the hopping-one convention. Scaling by $\lambda$ gives $|\cI'|=4(\lambda-1)$. Their erratum~\cite{AvMSerratum} concerns the continuum Mathieu example and leaves this discrete theorem intact.

For every $t\in[-R,R]$, Chambers realizes $P(E)=t$ as a Hermitian fibre spectrum. The $q$ monotone inverse branches therefore cover $[-R,R]$, with touching endpoints permitted. Removing $[-R+4,R-4]$ in their image leaves two pieces per branch. Change variables by $t=|P(E)|$. At a touching endpoint the integral is improper and integrable; endpoints have zero measure. This proves~\eqref{eq:surplus}.
\end{proof}

\begin{proposition}[Markov lower bound for the rational surplus]\label{prop:surplus}
For every reduced $p/q$,
\begin{equation}\label{eq:Markov}
 s_q\ge\frac{8(\lambda-1)}{q^2(\lambda^q+1)}
       \ge4(\lambda-1)q^{-2}e^{-Lq}.
\end{equation}
\end{proposition}
\begin{proof}
For a phase-union band $B$ of length $\ell$, Markov's polynomial inequality rescaled from $[-1,1]$ gives
\[
 \sup_B|P'|\le 2q^2R/\ell.
\]
The two end pieces in $B$ have total image length eight, so their total length is at least $4\ell/(q^2R)$. Sum over bands and use $|\cU|\ge4(\lambda-1)$ to obtain~\eqref{eq:Markov}.
\end{proof}

Avila--Last--Shamis--Zhou~\cite[Proposition~5.1]{ALSZ} already give quantitative rational band-end strip bounds with exponential rate $L$ after dual scaling. The surplus mechanism and that rate have a close published predecessor. Formula~\eqref{eq:Markov} makes the following deduction independent of the unpublished derivative estimates.

For $\eta=|\alpha-p/q|$, the window argument of Paper~A~\cite[Lemmas~2.5--2.6]{PaperA} gives
\begin{equation}\label{eq:continuity}
 d_{\mathrm H}(\cU(p/q),\Sig)\le\epsilon_q,
 \qquad \epsilon_q=(8+4\pi\lambda)\sqrt\eta+4\pi\lambda\eta.
\end{equation}
The irrational measure is $|\Sig|=4(\lambda-1)$ for every irrational frequency. For this noncritical statement one can cite Jitomirskaya--Krasovsky~\cite[Theorem~2]{JKmeasure}; Avila--Krikorian~\cite{AKmeasure} also cover the critical completion.

\begin{proposition}[Capped gap measure and a tail label]\label{prop:capped}
With $s_q$ and $\epsilon_q$ as above,
\begin{equation}\label{eq:capped}
 \sum_{m\ne0}\min\{|G_m|,2\epsilon_q\}\ge s_q-2\epsilon_q.
\end{equation}
If $s_q\ge16\epsilon_q$ and $M=\lfloor s_q/(8\epsilon_q)\rfloor$, some $|m|>M$ satisfies
\begin{equation}\label{eq:pigeonhole}
 |G_m|\ge\frac{s_qM}{8m^2}
       \ge\frac{s_q^2}{128\epsilon_qm^2}.
\end{equation}
\end{proposition}
\begin{proof}
The measure of $\cU\setminus\Sig$ is at least $s_q$. Its portions beyond the two spectral extremes have total length at most $2\epsilon_q$. Inside a bounded gap $(a,b)$, every point of $\cU$ lies in the $\epsilon_q$ endpoint strips. Their total length is at most $\min(b-a,2\epsilon_q)$. The bounded gaps have distinct integer labels. Countable additivity proves~\eqref{eq:capped}.

The $2M$ labels with $0<|m|\le M$ contribute at most $4M\epsilon_q\le s_q/2$. Since $2\epsilon_q\le s_q/8$, the uncapped tail has measure at least $s_q/4$. If all its gaps were less than $s_qM/(8m^2)$, summing the two signs and using $\sum_{r>M}r^{-2}\le1/M$ would contradict that bound. Finally $M\ge s_q/(16\epsilon_q)$.
\end{proof}

\begin{theorem}[The surplus range, including infinite $\beta$]\label{thm:measure}
If $\beta>2L$, then for every $2L<\beta'<\beta$ and $L<h<\beta'/2$, infinitely many labels satisfy
\[
 |G_m|\ge C|m|^{-K(h,\beta')},\qquad
 K(h,\beta')=2+\max\left\{0,\frac{4h-\beta'}{\beta'-2h}\right\},
\]
where $C>0$. In particular $\rho=0$ for $\beta=\infty$ as well.
\end{theorem}
\begin{proof}
On infinitely many scales $b_n\ge\beta'$,~\eqref{eq:continuity} gives $\epsilon_{q_n}\le C_1e^{-\beta'q_n/2}$. Equation~\eqref{eq:Markov} implies $s_{q_n}\ge c_he^{-hq_n}$. Thus $M_n\ge c e^{(\beta'/2-h)q_n}\to\infty$. The label supplied by~\eqref{eq:pigeonhole} obeys
\[
 |G_m|\ge c e^{(\beta'/2-2h)q_n}|m|^{-2},
 \qquad |m|\ge c e^{(\beta'/2-h)q_n}.
\]
If $\beta'<4h$, substitute the upper bound on $q_n$ from the second inequality into the negative exponential in the first. Otherwise discard the nonnegative exponential. This proves the formula. For $\beta=\infty$ choose any $\beta'>4h$ and obtain exponent two with a constant. The bounded spectral diameter proves $\rho\ge0$, as before.
\end{proof}

\begin{proposition}[Even denominators and small excess]\label{prop:even}
Let $\beta_e$ be the limsup of $b_n$ along indices with even $q_n$, with value zero if this set is finite. If $\beta_e>L$, then for every $L<\beta'<\beta_e$, $L<h<\beta'$ and $\epsilon>0$, infinitely many labels satisfy
\[
 |G_m|\ge C_\epsilon
 |m|^{-2-\max\{0,(2h-\beta')/(\beta'-h)\}-\epsilon}.
\]
The same conclusion, and hence $\rho=0$, holds with $\beta$ in place of $\beta_e$ if a $\beta$-attaining subsequence has a rational gap with excess
$\operatorname{ex}_n=|P_n(c_n)|-R_n=e^{o(q_n)}$.
\end{proposition}
\begin{proof}
For even $q$, the central rational closure gives $P(0)=\sigma R$, $P'(0)=0$, with $\sigma=(-1)^{q/2}$. The closure is classical (Bellissard--Simon~\cite{BellissardSimon1982}); the polynomial value is also given by Lamoureux--Mingo~\cite{LamoureuxMingo2007}. Paper~C~\cite[Corollary~7.2]{PaperC} fixes our normalization. The initial branch from zero to the point $E_4>0$ where $\sigma P(E_4)=R-4$ lies in $\cU$. The near-spectrum derivative bound of Paper~A, Proposition~3.3(2), applies on it eventually. If $K=\sup|P''|\le Cq^2e^{hq}$, then
\[
 4=R-\sigma P(E_4)\le KE_4^2/2,\qquad
 s_q\ge2E_4\ge cq^{-1}e^{-hq/2}.
\]
Apply Proposition~\ref{prop:capped} with $\epsilon_q\le C e^{-\beta'q/2}$. This gives
\[
 M\ge cq^{-1}e^{(\beta'-h)q/2},\qquad
 |G_m|\ge cq^{-2}e^{(\beta'/2-h)q}|m|^{-2}.
\]
The same power conversion as above gives the assertion; powers of $q=O(\log|m|)$ are absorbed by the arbitrary $\epsilon$.

For the excess variant first put the critical point in the derivative domain. If the rational gap has length $w$ and diameter bound $D_0=4+4\lambda$, real-root factorization gives
\[
 (\log|P|)''=-\sum_j(E-e_j)^{-2}\le-q/D_0^2.
\]
Integrating to the more distant endpoint shows
\[
 \frac{qw^2}{8D_0^2}\le\log(1+\operatorname{ex}/R)
 \le\operatorname{ex}/R.
\]
Thus $w\to0$ exponentially when $\operatorname{ex}=e^{o(q)}$. The gap and its critical point are near $\Sig$ by~\eqref{eq:continuity}. On an adjacent end piece the bound $|P''|\le K$ is therefore valid. If $v=|P'(x)|$, reaching $P'(c)=0$ takes distance at least $v/K$. Integrating $v-Kt$ along that segment gives
\[
 |P'(x)|^2\le2K(|P(c)|-|P(x)|)
          \le2K(\operatorname{ex}+4).
\]
The inverse-derivative integral then yields $s_q\ge4/\sqrt{2K(\operatorname{ex}+4)}=e^{-(h/2+o(1))q}$. The previous argument applies on the stipulated same resonant subsequence.
\end{proof}

These are sufficient conditions. No converse estimate on the actual surplus is established here. Nor does this proof identify the individual labels of rational campaign gaps with their minimal rational representatives at an irrational frequency. If a rational gap has width $w>2\epsilon_q$, its central core lies in an irrational gap of width between $w-2\epsilon_q$ and $w+2\epsilon_q$. This is an absolute width statement about the containing gap. It does not, by itself, identify its exact integer label.

\label{end:surplus}

\section{The arithmetic picture and certified finite checks}\label{sec:diagram}
\subsection{A phase diagram with phase conditions}
Put $x=2\beta/L$. Denote the upper Abel-averaged fastest exponent at phase $\theta$ by $\alpha_u^+(\theta)$ and its supremum over phases by $\alpha_{\rm worst}$. Paper~C also defines an exponent uniform over launch sites, $\alpha_u^{\rm unif}$, which bounds that supremum. These quantities agree in the exact-law regime specified below.

The typical fastest exponent and the two packing dimensions of every nonzero spectral measure are $\max\{0,1-2/x\}$ at almost every phase, by Paper~A~\cite[Theorems~A and C]{PaperA}; the value is zero at $x=0$. The packing entries concern spectral measures, not the packing dimension of the spectral set. At $\beta=\infty$ these values are one at every phase. Paper~A gives the finite-positive-$\beta$ law used here. Its $\beta=\infty$ transport and packing endpoint is due to Jitomirskaya--Zhang~\cite[Corollary~3]{JZ2022}, and its $\beta=0$ fastest endpoint to Damanik--Tcheremchantsev~\cite{DT2007}. Paper~A sharpens earlier packing estimates of Jitomirskaya--Liu--Tcheremchantsev~\cite{JLT2025}. The typical transport framework for general supercritical type~I operators is developed in Paper~B~\cite{PaperB}.

Table~\ref{tab:phase} keeps the phase quantifiers visible. ``AL'' means Anderson localization and ``SC'' purely singular continuous spectrum. Completely resonant phases satisfy $2\theta\in\alpha\Z+\Z$. The phase $1/4$ is not completely resonant for irrational $\alpha$.

\begin{table}[ht]
\centering\footnotesize
\setlength{\tabcolsep}{4pt}
\begin{tabular}{@{}>{\raggedright\arraybackslash}p{.08\linewidth}>{\raggedright\arraybackslash}p{.23\linewidth}>{\raggedright\arraybackslash}p{.26\linewidth}>{\raggedright\arraybackslash}p{.23\linewidth}>{\raggedright\arraybackslash}p{.12\linewidth}@{}}
\toprule
$x$ & Spectral type at a.e. phase & Specified exceptional phases & $\alpha_{\rm worst}$ & $\rho$\\
\midrule
$0$ & AL & AL at completely resonant phases; SC at a dense set of other phases & $0$ & $1$ if $\alpha$ is Diophantine\\[3pt]
$(0,1)$ & AL & AL at completely resonant phases & $[0,x/(x+2)]$ & $[0,1-x]$; $\rho\ge1/C$ if $x\le2/C$\\[3pt]
$1$ & AL & No equality-line classification from the cited strict criterion & $[0,1/3]$ & Open\\[3pt]
$(1,2)$ & AL & SC at $1/4$ if $\beta_e=\beta$ & $[1-1/x,x/(x+2)]$ if $\beta_e=\beta$ & $0$\\[3pt]
$2$ & Not classified by the strict arithmetic transition & SC at $1/4$ if $\beta_e=\beta$ & $1/2$ if $\beta_e=\beta$ & $0$\\[3pt]
$(2,\infty)$ & SC at every phase & Same all-phase statement & $1-1/x$ if $\beta_e=\beta$ & $0$\\[3pt]
$\infty$ & SC at every phase & Same all-phase statement & $1$ & $0$\\
\bottomrule
\end{tabular}
\caption{Established statements and explicit open entries. Every bracket in the transport column is a pair of proved bounds. The gap column uses Theorems~\ref{thm:below}, \ref{thm:polynomial} and~\ref{thm:measure}, and Corollary~\ref{cor:bracket}. No exact gap law is asserted at $x=1$.}\label{tab:phase}
\end{table}

The a.e.-phase AL and all-phase SC statements use Avila--You--Zhou~\cite[Theorem~1.1]{AYZ2017} for $0<\beta<\infty$, with strict inequalities at $L=\beta$. The endpoint spectral-type assertions at $x=0$ and $x=\infty$ use the classical endpoint results recalled in Avila--You--Zhou (Remark~1.2 of arXiv:1512.03124). Liu~\cite[Theorem~1.1]{Liu2022} proves completely resonant AL for $L>2\beta$. At $1/4$, even $q_n$ imply a phase resonance at $q_n/2$ with error $|\delta_n|/2$; its exponent is at least $2\beta_e$. The absence-of-point-spectrum part of Jitomirskaya--Liu~\cite[Theorem~4.2 and the remark after Theorem~1.1]{JLphase} therefore gives SC when $L<2\beta_e$. Its localization half has additional frequency hypotheses; we do not use that half for this entry.

The transport bounds are from Paper~C, Theorem~7.5 and Corollaries~5.6 and~7.6. The phase-$1/4$ lower bound uses even denominators and $\beta_e$, not even indices. Without $\beta_e=\beta$, retain that lower bound with $\beta_e$ and the typical all-phase lower bound separately. The exact worst-phase formula is then not supplied by these imports. A conditional parity-free mechanism in Paper~C requires its own wide-band hypothesis; it is not a consequence of the gap theorem here.

At a finite scale,
\begin{equation}\label{eq:finitecross}
 \frac{\log\kappa_n}{Lq_n}
 =1-\frac{2b_n}{L}+\frac{2\log(2\pi q_n)}{Lq_n}.
\end{equation}
Thus $x=1$ does not force $\kappa_n=O(1)$. Fixed margins on either side of one give exponential smallness or largeness along the corresponding scales. The logarithmic correction matters at the equality line.

\subsection{The known transport geometry}
For even $q$, put $\sigma=(-1)^{q/2}$ and let $0<\eta\le1/4$. In one period of the local phase drift, the launch set
\[
 \left|\cos2\pi(q\theta+j\delta)-\sigma\right|
                     \le\eta\Lambda^{-1}/16
\]
has cardinality $q\sqrt{\eta/(2\kappa)}(1+O(\Lambda^{-1}))+O(1)$. This is the quadratic turning-point width divided by the mesh $|\delta|$. Its logarithm has rate $(b-L/2)q$ only when the main count diverges. Continuous window length is not an exact test for occupancy by one lattice cell.

Paper~C's nearby-phase burst theorem applies with its scale condition $b\ge h/2+(h-L)$ and time interval $e^{(h/2+(h-L))q}\le T\le|\delta|^{-1}$. It gives a polynomially small fastest tail. At the resonant time its radius has exponent $b-h/2$, approaching the preceding geometric exponent as $h\downarrow L$.

The confinement estimate of Paper~C, Proposition~6.4, is a uniform square-root upper bound on a cosine sublevel arc, plus its collars. At a generic phase a particular arc can be narrower. In Paper~A's notation $J_n=3+C_hq^{5/2}e^{hq}|\delta|$ is greater than three; it is $3+o(1)$ on $b\ge L+\gamma$ if $h-L<\gamma$. Alternatively $\|S^q+S^{-q}+R_q\|\le3$ eventually there. These bounds match the turning-point exponential radius. Below $b=L$, the coarse norm bound permits a larger confinement upper bound; it does not prove that this enlargement is attained by the dynamics.

\subsection{What the numerical checks certify}\label{sec:numerics}
All numbers in this subsection are inherited from the independent verification reports D77 and D80. They are interval certificates or exact finite-matrix counts. We do not rerun those campaigns here. No number is used in a proof.

For the Mathieu gaps $G_D(\kappa)$ of Appendix~\ref{app:mathieu}, D77 encloses a crossing $G_1=2G_2$ in
\[
 \kappa_*\in[1.4260784738\mathbin{\pm}6.60\cdot10^{-11}].
\]
This certifies a crossing, not global uniqueness. Ten independent checks of the rational phase-union prediction give the intervals in Table~\ref{tab:collapse}. Here the parent is $p/q$, $A$ is the following continued-fraction coefficient in the verified family, $w$ is the measured rational width, and $P'$ is the parent derivative at the measured location. This tests widths conditional on locations.

\begin{table}[ht]
\centering\small
\begin{tabular}{@{}rrrrl@{}}
\toprule
Parent & $\lambda$ & $A$ & $|m|$ & Enclosure of $\log(w|P'|/G_D)$\\
\midrule
$2/3$ &5&699&3&$[0.0001979253,0.0001979254]$\\
$3/5$ &5&2000&10&$[0.0000162096,0.0000162097]$\\
$3/5$ &2&600&15&$[-0.0005267205,-0.0005267204]$\\
$8/13$&5&385&78&$[0.0039697174,0.0039697175]$\\
$8/13$&2&77&52&$[0.0172042988,0.0172042989]$\\
$5/8$ &2&200&16&$[0.0003463309,0.0003463310]$\\
$3/5$ &5&60&15&$[0.0091106244,0.0091106245]$\\
$2/3$ &5&69&3&$[0.0034522615,0.0034522616]$\\
$8/13$&2&385&13&$[0.0005374229,0.0005374230]$\\
$2/3$ &5&149&6&$[0.0008759784,0.0008759785]$\\
\bottomrule
\end{tabular}
\caption{Certified finite rational Mathieu collapse from D77, Section~5. There are ten cases and eight distinct integer labels. These data do not prove an irrational gap law.}\label{tab:collapse}
\end{table}

D80 tests the translated construction itself on boxes containing several wells. All three cases have parent $5/13$. Case G uses $\lambda=2$, $\theta=0.0377$, $\hbar=0.1$ and three wells. H uses the same coupling and phase, $\hbar=0.005$ and two wells. OOS was registered before its independent measurement: $\lambda=3$, $\theta=0.0193$, $\hbar=0.1$ and three wells. The arithmetic uses $\delta=\hbar/(2\pi q\lambda^{q/2})$.

\begin{table}[ht]
\centering\small
\begin{tabular}{@{}llcl@{}}
\toprule
Case & Inclusive box & Counts for $D=1,\ldots,6$ & Maximum level error\\
\midrule
G & $[731,222520]$ & $3,6,9,12,15,18$ & $[0,0.000282844609]$\\
H & $[14638,2971820]$ & $2,4,6,8,10,12$ & $[0,0.000314943324]$\\
OOS & $[256912,3351000]$ & $3,6,9,12,15,18$ & $[0,0.000013630460]$\\
\bottomrule
\end{tabular}
\caption{D80 finite-box certificates. Errors are in units of $2\hbar/P_*'$ after the common dipole shift. Each row includes three interior fractions of each counting window and midpoint counts after moving both ends by $-1000$ or $+1000$ sites.}\label{tab:globalchecks}
\end{table}

The ground-cluster projector assigns each translated anchor class the per-well share
\[
\begin{split}
 \text{G: }&[0.9558209196836182962179037231977
                 \mathbin{\pm}1.93\cdot10^{-32}],\\
 \text{OOS: }&[0.98136338757751036956774744137
                 \mathbin{\pm}3.53\cdot10^{-30}].
\end{split}
\]
The tables and the display above report the independent certificates, including the corrected cluster-projector weights. The original single-eigenvector weight prediction, residual-growth prediction and factor-30 switch prediction for the ratio of a gap to the derivative proxy TP1 failed; TP1 is not a measured transport exponent. At $q=21$ and $\lambda=5$, the literal TP1 interval also failed, and the full original gap-switch calculation remains uncertified. These finite checks do not establish asymptotic complementarity. The accompanying Numerical Record, \S\S1--3, preserves the registrations, their certified comparison intervals, the corrected projector interpretation and the still uncomputed full-period test.

\subsection{New finite certificates below the threshold}\label{sec:below-numerics}
The three new independent checking campaigns certify finite rational widths, finite rounded-span Green columns, and finite rounded-span complementary spectra. These are different experiments. They do not certify exact infinite-support Wannier functions or an irrational asymptotic theorem. The accompanying Numerical Record, \S\S4--6, gives the fuller intervals and preserves failed registrations. No new computation was used to produce the intervals printed here.

The rational-width campaign independently certified 19 of 64 submitted widths and four registered out-of-sample widths. The latter use $\lambda=5$, parent $p/q=13/21$, and extensions
$(13A+8)/(21A+13)$. Their certified fibre counts are $D$ above the centre, fixing rational IDS $1-D/(21A+13)$.
\begin{table}[ht]
\centering\small
\begin{tabular}{@{}rrl@{}}
\toprule
$A$&$D$&Outward enclosure of $\log(\text{rational gap width})$\\
\midrule
20&2&$[-68.974899318838,-68.974899318837]$\\
20&3&$[-101.720353379131,-101.720353379130]$\\
40&2&$[-67.646792273466,-67.646792273465]$\\
40&3&$[-99.079273854345,-99.079273854344]$\\
\bottomrule
\end{tabular}
\caption{Finite rational certificates at the two registered $D$ values. They are not irrational-transfer certificates.}\label{tab:belowwidthchecks}
\end{table}

For the continuation experiment define the finite-span statistic
\[
 \rho_c=\frac{E_*-\sup\sigma(Q_IH_IQ_I|_{\Ran Q_I})}{m_0}.
\]
The exact dyadic frame coefficients define the tested span.
At $\lambda=5$, parent $13/21$, offset $y=3/10$, the independently certified enclosures are
\[
\begin{array}{ccl}
 b&\text{family}&\rho_c\\
 .28&\text{continued and box}&[66.310252694963,66.323516071840]\\
 .33&\text{continued and box}&[49.271919945713,49.281775315240].
\end{array}
\]
Each family was certified separately; the repeated enclosure does not assert exact equality of their spectra. Every case also passed the strict test $\rho_c>1.000001$. Here the numerical clamp is the specified $C^1$ cubic clamp, whereas the analytic construction uses a fixed $C^\infty$ clamp.

The Green-column campaign independently certified six original rounded-span boxes and two new ones. It measures a cell-boundary tail-rate statistic, rather than the full infinite-volume Schur norm. Its exact window convention and all eight outward intervals are in Numerical Record \S5.

The rational $D=5$ slope registration within $0.25$ of $2(D-1)$ failed in its tested family. The universal corrected-increment residual bound $0.0061$ was also too small; a certified value has magnitude greater than $0.006597244239$. The safe original-sample bound is $0.00660$. In the continuation campaign the original universal no-twofold-drop and finite-$q$ continued/box agreement predictions failed in some preasymptotic cases. The later registered out-of-sample bands passed. These successes do not retroactively change failed registrations and no finite sample proves a limiting constant.

\label{end:diagram}

\section{Conjectures}\label{sec:conjectures}
This section contains no further theorem. The global gap rate in~\eqref{eq:rho} must be distinguished from the ratio between adjacent satellites at one rational scale.

\begin{conjecture}[The unresolved gap side]\label{conj:gap}
For $0<x=2\beta(\alpha)/L<1$,
\[
 \rho(\lambda,\alpha)=1-x.
\]
\end{conjecture}
The perturbative Mathieu formula at fixed $D$ suggests this value. It does not prove it. For $\kappa\gg1$, the $D$-th model gap is of order $\kappa^{1-D}$, much smaller than an error $o(\sqrt\kappa)$ when $D\ge2$. Resolving these gaps needs relative tunnelling estimates. An exact global rate also needs control over all sufficiently large labels. A local ladder supplies neither requirement.

\begin{conjecture}[The missing dynamical upper bound]\label{conj:transport}
If $1<x<2$, then
\[
 \alpha_u^{\rm unif}\le1-1/x.
\]
Together with the known lower bound this would give the exact worst exponent when $\beta_e=\beta$. A parity-free exact formula additionally requires a parity-free lower mechanism.
\end{conjecture}
A band-edge spectral comparison alone does not establish this conjecture. Position commutators, launch-state estimates, complementary-energy bounds and control between resonant scales are also needed.

\begin{conjecture}[Complementarity]\label{conj:complementarity}
The gap and worst-transport arithmetic laws are
\[
 \rho=(1-x)_+,\qquad
 \alpha_u^{\rm unif}=\alpha_{\rm worst}=(1-1/x)_+,
\]
with the appropriate endpoint interpretations. Equivalently, the proposed functions satisfy $\alpha_{\rm worst}(x)=\rho(1/x)$.
\end{conjecture}
The vanishing gap half for $x>1$ is now a theorem. The other equalities remain open in the ranges described above. Below one, zero fastest exponent is a weaker conjecture than dynamical localization. Literal all-phase dynamical localization is false: Jitomirskaya--Simon~\cite{JSimon} give dense sets of phases with no point spectrum. In the supercritical regime these have SC spectrum, precluding dynamical localization.

A stronger quantitative ladder beyond the second cap of~\eqref{eq:Dcap} is another open problem. A proof would have to retain both the inverse-polynomial function and the cosine potential inside a uniform quantization theorem. We do not import a general perturbation-stable Harper theorem under the name of the Mathieu ladder. Such a theorem would have additional parameter-uniformity obligations.

\label{end:conjectures}

\section{Discussion}\label{sec:discussion}
The proof begins with the highest rational band, where convexity and Chambers' formula provide a contour with a large discriminant. The stable Floquet line then gives localization per physical site. That localization is what makes translations, Gram summation and hard frame cutoffs controllable as the denominator grows.

The orbital term moves the common energy origin of the wells. It does not change the count between consecutive clusters. Its Wannier form is a Brillouin-zone average; it is not, in general, its value at Bloch momentum zero. The off-diagonal first-order terms are suppressed by tunnelling. The remaining comparison has two independent relative errors. The two caps on $D$ control those errors separately. Keeping them separate is essential near $b=L/2$.

The final integer count is the bridge from a local well calculation to an irrational gap. Correct covariance gives one symbol in every well. The common Mathieu windows give the same number $D$ in both comparison operators. Their equality forces an exact count $KD$ in each box. No inference from two agreeing numerical box spectra is used in this passage.

The theorem gives a lower bound along a sequence of labels. The endpoint $\beta=L/2$, bounds for all large labels and a matching dynamical law remain open. The change of exponent at $3L/2$ records the precision of the quadratic comparison.

Below the threshold the relative tunnelling scale replaces the small-curvature Mathieu spacing. The continuation is periodic at both far ends and retains only the selected physical well. Its exact covariance makes the full diagonal constant common to all offsets. Keeping the self-energy and estimating its first phase derivative is sufficient for the quadratic model hypothesis; no cubic Taylor assertion across frame jumps is needed.

\subsection*{How this work was done}
The authors developed the results and proofs with assistance from AI models from more than one developer. They take responsibility for every claim. The authors chose the question, directed the work and decided which lines to pursue; the models did most of the exploration, derivation and drafting. Both authors contributed equally.

The source work proceeded through separate research sessions on gap sizes, the arithmetic picture, the polynomial remainder, top-band analyticity and the global reduction. Each was followed by a separate session instructed to refute the submitted argument and inspect its primary sources; these checks are D76--D80 in the working archive, and their repairs take precedence. The initial global-frame argument failed its check; the translated-well construction presented here replaced it. A further refutation-first audit of the assembled manuscript (D81) required repairs to Appendix~\ref{app:mathieu}, to the companion table and to several attributions; this version makes them.

The new below-threshold gap argument, the complementary Green-function estimate and the corrected covariant continuation were developed in separate research sessions, each followed by an independent refutation-first check whose replacement estimates are incorporated here.

Priority and literature searches were AI-assisted; an unsuccessful bounded search is not a priority certificate. No computation is used in any proof.

\label{end:discussion}

\section{Methods and conventions}\label{sec:methods}
We use the physical shift $(Su)(j)=u(j-1)$, so $S^c\delta_0=\delta_c$, and $X\delta_j=j\delta_j$. Reflection is $(\mathcal Ru)(j)=u(-j)$. The potential is $v(t)=2\lambda\cos2\pi t$. All inner products are conjugate-linear in the first argument. Analytic continuations of matrix entries are taken before any conjugation used on the real axis.

Chambers' formula, in the uniform Bloch gauge, is
\begin{equation}\label{eq:Chambers}
 \det(E-B_t(k))=P(E)-2\cos k-c(t),\qquad
 c(t)=2\Lambda\cos2\pi qt.
\end{equation}
Here $B_t(k)$ has cyclic hopping factors $e^{\pm ik/q}$. The sign of $k$ is fixed so that $[A,X](k)=-iq\partial_kA(k)$. The physical shift or its adjoint can be used as the cyclic hopping matrix; this declaration fixes the orientation in the orbital formulas. All energy spectra and all width estimates are invariant under reversing $k$.

Six inputs organize the proof. (I1) is Chambers' identity. (I2) is the near-spectrum derivative estimate in Paper~A, Proposition~3.3(2). (I3) is the complex-energy extension of Paper~B's growth and rational-product lemmas, proved in Appendix~\ref{app:complex}. (I4) is the Aubry--Perron maximal-site fact at $B_0(0)$, proved in Appendix~\ref{app:frames}. (I5) is the common uniform Mathieu ladder, derived from one-well quantization in Appendix~\ref{app:mathieu}. (I6) comprises the continuous IDS, full DOS support, and the convention $N=\{m\alpha\}$; see Avron--Simon~\cite{AvronSimon} and Avila--You--Zhou~\cite[Section~1]{AYZdry}.

For fixed $\xi\in(0,L)$ we set $\tau=\xi/100$. Only finitely many derivatives and resolvent products occur. Let $C_*$ dominate their exponential-loss coefficients. We choose $h>L$ only after that list is fixed, so that
\begin{equation}\label{eq:slack}
 C_*(h-L)<\min\{\xi,\tau,L\}/10000.
\end{equation}
The notation $\EE=e^{o(q)}$ in the appendices denotes a positive quantity bounded by $Cq^Ce^{C(h-L)q}$ for one of this finite list of constants. Each occurrence can be bounded by an arbitrarily small fixed exponential loss by this choice and a subsequent large-$q$ threshold. It never hides dependence on $b,D$, the well or $\theta$. To obtain a literal $o(1)$ in an exponent, diagonalize the fixed-loss choices and their thresholds. The extension of Theorem~\ref{thm:scale} to $\xi=L$, if desired, follows by applying any smaller margin and restricting to its single permitted label $D=1$.

Dirichlet compression to an interval $I$ is denoted by a subscript $I$. Eigenvalues are counted with multiplicity. Strict inequalities in counting windows are used only away from the enclosing clusters. Schur norms always include both row and column sums. Weighted Schur estimates are proved with their spatial weights before the corresponding norm is simplified.

\label{end:methods}

\appendix
\section{Complex energies and rational transfer products}\label{app:complex}
This appendix proves the extension of Paper~B, Lemmas~2.4--2.5, used below. It does not assert that the complex-energy Lyapunov exponent equals its real spectral value.

For $z\in\C$, put
\[
 A(z,t)=\begin{pmatrix}z-v(t)&-1\\1&0\end{pmatrix},\qquad
 A_r^{(\omega)}(z,t)=A(z,t+(r-1)\omega)\cdots A(z,t).
\]
On the spectrum of the supercritical almost Mathieu operator, the Lyapunov exponent equals $L=\log\lambda$, by Bourgain--Jitomirskaya~\cite[Corollary~2]{BourgainJitomirskaya2002}, after matching the coupling convention. This spectral identity is the AMO instance of the growth input in the companion papers.

\begin{lemma}[Uniform complex neighbourhood]\label{lem:complexgrowth}
For each $h>L$ there are a bounded open neighbourhood $U_h^{\C}$ of $\Sig$ and $C_h<\infty$ such that
\[
 \norm{A_r^{(\alpha)}(z,t)}\le C_he^{hr}
 \quad(z\in U_h^{\C},\ t\in\T,\ r\ge0).
\]
After enlarging $C_h$, every sufficiently large convergent $p/q$ satisfies
\begin{equation}\label{eq:complexrat}
 \norm{A_r^{(p/q)}(z,t)}\le C_he^{hr}
 \quad(0\le r\le q, z\in U_h^{\C},\ t\in\T).
\end{equation}
The same estimates hold for subblocks.
\end{lemma}
\begin{proof}
Choose $h_1\in(L,h)$. At each $E_0\in\Sig$, the infimum formula for the Lyapunov exponent gives an integer $r_0$ with
\[
 \frac1{r_0}\int_\T\log\norm{A_{r_0}^{(\alpha)}(E_0,t)}\,dt<h_1.
\]
The finite product is jointly continuous in complex $z$ and real $t$. Its determinant is one, so its norm never vanishes. The strict averaged inequality persists on a compact complex neighbourhood $K_0$ of $E_0$.

Rotation by $r_0\alpha$ is uniquely ergodic. The continuous functions $t\mapsto\log\norm{A_{r_0}^{(\alpha)}(z,t)}$, $z\in K_0$, form a compact family in the uniform norm. A finite uniform approximation to that family makes the Birkhoff convergence uniform in both variables. For all sufficiently many blocks, their summed logarithmic norms are at most $mr_0h$. Splitting $r=mr_0+s$ leaves a uniformly bounded remainder. This proves the bound on $K_0$. A finite cover of $\Sig$ gives $U_h^{\C}$ and one constant.

For the rational comparison, each one-step error has norm at most $d=\Lip(v)|\delta|$. A term with $j$ such errors in a product of $r\le q$ factors contains at most $j+1$ irrational subblocks. Its total norm is bounded by $C_h^{j+1}e^{h(r-j)}d^j$. Consequently
\[
 \norm{A_r^{(p/q)}}\le C_he^{hr}(1+C_hde^{-h})^r
 \le C_he^{hr}\exp(C_h\Lip(v)e^{-h}),
\]
because $q|\delta|<1$ for convergents. No step uses realness of $z$. Starting phases cover all subblocks.
\end{proof}

\begin{corollary}[Polynomial derivatives and the top contour]\label{cor:derivatives}
For fixed $r\ge1$, on a fixed smaller complex neighbourhood,
\[
 |\partial_z^r\Delta_t(z)|\le C_{r,h}q^re^{hq},\qquad
 |P^{(r)}(z)|\le C_{r,h}q^re^{hq}.
\]
A fixed-radius complex neighbourhood of $E_*(p_n/q_n)$ lies in $U_h^{\C}$ for all sufficiently large $n$.
\end{corollary}
\begin{proof}
Differentiate the product. Each derivative inserts $\left(\begin{smallmatrix}1&0\\0&0\end{smallmatrix}\right)$ at a distinct position. There are at most $q^r$ ordered choices and at most $r+1$ rational segments. Apply~\eqref{eq:complexrat} and bound the trace by twice the norm. The Chambers correction is independent of energy. Finally~\eqref{eq:continuity} gives convergence of the rational top edge to $\max\Sig$; compact containment in the open complex neighbourhood supplies the fixed radius.
\end{proof}

For real energies this is also exactly the derivative input (I2) of Paper~A, Proposition~3.3(2). Before using it on a whole rational band,~\eqref{eq:continuity} puts the whole phase union in its neighbourhood. It is never applied at an unanchored critical point far from the spectrum.

\label{end:complex}

\section{The rational top band, Green decay and flatness}\label{app:top}
\subsection{Convexity and an isolated inverse branch}
Fix a sufficiently large convergent and put $\Lambda=\lambda^q$, $R=2\Lambda+2$. Let $I_q=[a_q,E_*]$ be the top phase-union band and let $P_*'=P'(E_*)$.

\begin{lemma}[Top-band geometry]\label{lem:convex}
The polynomial is convex on the final monotonicity interval, and
\begin{equation}\label{eq:topbounds}
 P_*'=\max_{I_q}P'\ge\frac{2R}{|I_q|}
       \ge\frac{\Lambda+1}{\lambda+1},\qquad
 |I_q|\ge c_hq^{-1}e^{-(h-L)q}.
\end{equation}
All roots of $P-R$ other than $E_*$ lie at distance at least $|I_q|$ from $E_*$.
\end{lemma}
\begin{proof}
For each $w\in[-R,R]$, Chambers' formula realizes $P-w$ as the characteristic polynomial of a Hermitian fibre. Rolle's theorem gives real critical points with critical modulus at least $R$. The zeros of $P''$ interlace those of $P'$. Hence $P''>0$ to the right of the last critical point, which is at or below $a_q$. Integrating $P'$ on $I_q$ gives $2R\le |I_q|P_*'$. The spectral diameter is at most $4+4\lambda$, proving the second inequality. The derivative upper bound in Corollary~\ref{cor:derivatives} proves the band-length lower bound. The other roots belong to lower monotone bands, hence are no larger than $a_q$.
\end{proof}

Set
\begin{equation}\label{eq:contour}
 r=\frac{|I_q|}{4q},\qquad \mathcal C_r=\{z:|z-E_*|=r\}.
\end{equation}
Initially assume $0\le2\Lambda-c(t)\le1$.

\begin{lemma}[Contour and top inverse]\label{lem:contour}
On $\mathcal C_r$,
\[
 |P(z)-R|\ge\tfrac34rP_*'\ge\frac{3R}{8q},\qquad
 |\Delta_t(z)|\ge\frac{\Lambda}{2q}.
\]
For $|w-R|\le R/(8q)$, $P(E)=w$ has exactly one root in $|E-E_*|<r$, and this root lies in $|E-E_*|<r/2$. It is simple and analytic in $w$.
\end{lemma}
\begin{proof}
Factor $P(z)-R=(z-E_*)\prod_{j<q}(z-e_j)$. On either radius $r$ or $r/2$, each other factor is at least its value at $E_*$ times $1-1/(4q)$. Their product is at least three quarters of its value there. The radius-$r$ bound follows. On radius $r/2$ the bound is at least $3R/(16q)>R/(8q)$. Rouch\'e's theorem gives one root counting algebraic multiplicity on each circle. That root is therefore simple. Since $|R-c(t)|\le3$, subtraction gives the discriminant estimate eventually.
\end{proof}

\subsection{The decay is per physical site}
Let $G_t(x,y;z)=\inner{\delta_x}{(B_t-z)^{-1}\delta_y}$.

\begin{lemma}[Periodic Green function]\label{lem:green}
For $z\in\mathcal C_r$,
\begin{equation}\label{eq:green}
 |G_t(x,y;z)|\le K_qe^{-L_q|x-y|},\qquad
 K_q=128C_h^2q^2e^{2(h-L)q},\quad
 L_q=L-\frac{\log(8q)}q.
\end{equation}
\end{lemma}
\begin{proof}
The cell monodromy has eigenvalues $\mu,\mu^{-1}$ with
$|\mu|\ge|\Delta_t(z)|-1\ge\Lambda/(4q)>1$. Let $\Psi_+$ span its stable line and $\Psi_-$ its unstable line. The reciprocal sine of their Hermitian angle is the norm of the unstable spectral projection. The formula
\[
 \Pi_u=\frac{M_x-\mu^{-1}}{\mu-\mu^{-1}}
\]
and~\eqref{eq:complexrat} give
\[
 (\sin\angle(\Psi_+,\Psi_-))^{-1}
 \le8qC_he^{(h-L)q}.
\]
For $0\le j<q$, propagate the stable solution backwards from $x+q$. A determinant-one $2\times2$ matrix has inverse with the same operator norm. Thus
\[
 \frac{\norm{\Psi_+(x+j)}}{\norm{\Psi_+(x)}}
 \le C_he^{h(q-j)}|\mu|^{-1}
 \le4qC_he^{(h-L)q}e^{-hj}.
\]
The Green-function Wronskian formula bounds its kernel by this quotient times the reciprocal sine. Every further complete cell contributes $|\mu|^{-1}$. Writing $y-x=mq+j$ gives~\eqref{eq:green} for $x\le y$, with the displayed enlargement of constants. For complex energy use transpose symmetry $G_t(x,y;z)=G_t(y,x;z)$ to cover $x>y$. Hermitian symmetry at complex $z$ is not asserted.
\end{proof}

The rate in~\eqref{eq:green} measures actual site distance, including the displacement inside a cell. No ordering of the sampled potential values is assumed.

\subsection{Complex Bloch strips}
Fix $\tau\in(0,L)$ and set $\eta_q=(L_q-\tau)q$.

\begin{lemma}[Analytic fibre data]\label{lem:strip}
For $|\operatorname{Im}k|\le\eta_q$ and $z\in\mathcal C_r$, the fibre resolvent exists and
\[
 \norm{(B_t(k)-z)^{-1}}\le K_q':=\frac{2K_q}{1-e^{-\tau}}.
\]
There is exactly one simple eigenvalue $E_t(k)$ inside $\mathcal C_r$, with $|E_t(k)-E_*|<r/2$. It satisfies
$P(E_t(k))=2\cos k+c(t)$. Its Riesz projection $\Pi_t(k)$ and reduced resolvent
$\mathcal S_t(k)=Q_t(k)(B_t(k)-E_t(k))^{-1}Q_t(k)$ are analytic, with
\begin{equation}\label{eq:Rieszbounds}
 \norm{\Pi_t(k)}\le rK_q',\qquad
 \norm{\mathcal S_t(k)}\le2K_q',\qquad
 \norm{\partial_kB_t(k)}\le2e^L/q.
\end{equation}
\end{lemma}
\begin{proof}
The uniform-gauge fibre resolvent is the sum of full-line kernels with weight $e^{ikd/q}$, where $d$ is the physical displacement. Its absolute value is bounded by $K_qe^{-\tau|d|}$. Both Schur sums are bounded by $K_q'$. Absolute convergence permits multiplication by $B_t(k)-z$ and proves the resolvent identity. Chambers and Lemma~\ref{lem:contour} select one root since $|2\cos k+c(t)-R|\le R/(8q)$ for large $q$. Riesz calculus gives the first two estimates: the contour has length $2\pi r$ and $|z-E_t(k)|\ge r/2$. The hopping derivative gives the last.
\end{proof}

All these estimates extend to
\[
 \Phi_{\rm fr}=\{t:0\le2\Lambda-c(t)\le M'\},\qquad M'=\Lambda e^{-\xi q/2},
\]
after replacing $\eta_q$ by $\eta' =\eta_q-\log4$. Indeed $M'=o(R/q)$, which preserves the contour inequalities. The scalar estimates apply on every phase component. A frame anchored at zero will be used only on the component about zero.

\begin{proposition}[Frozen orbital flatness]\label{prop:flatness}
At a simple top-band point write $D=\partial_tB$, $K=\partial_kB$ and $\mathcal S=\mathcal S_t(k)$. The top diagonal entry of the first-order polynomial symbol in Appendix~\ref{app:remainder} is
\begin{equation}\label{eq:traceflat}
 \mathcal G_{tt}(k)=\frac{qP'(E_t(k))}{2i}
 \bigl(\tr(\Pi D\mathcal SK)-\tr(\Pi K\mathcal SD)\bigr).
\end{equation}
Define $S_{\rm top}=-\mathcal G_{tt}/P'(E_t(k))$. For every fixed $j\ge0$ and $\varepsilon>0$, eventually
\[
 \sup_{k\in\R}\left|
 \partial_k^j\bigl(S_{\rm top}(k)-S_{\rm top}(0)\bigr)\right|
 \le C_{j,\varepsilon}e^{-(L-\varepsilon)q}.
\]
For $j\ge1$ the derivative of the constant is read as zero.
\end{proposition}
\begin{proof}
On the real axis the orbital formula is
$\mathcal G_{tt}=qP'(E_t)\operatorname{Im}\inner{u}{D\mathcal SKu}$. Taking the difference with its conjugate gives~\eqref{eq:traceflat}. The right-hand side is its analytic continuation. For a rank-one idempotent, even when it is not orthogonal, $|\tr(\Pi A)|\le\norm\Pi\norm A$. Equation~\eqref{eq:Rieszbounds} bounds the trace expression after cancellation of $P'$ by $\EE$. It is periodic under the uniform-gauge covariance. Its Fourier coefficients are therefore bounded by $\EE e^{-\eta'|m|}$. Sum the nonconstant coefficients for $j=0$ and sum $|m|^je^{-\eta'|m|}$ for each fixed $j\ge1$. Choose $h-L$ and the strip slack arbitrarily small to obtain the asserted rate. This uses the Fourier series directly and does not lose half the strip through a Cauchy derivative bound.
\end{proof}

The normalization has no extra factor $q$ in $S_{\rm top}=-\mathcal G_{tt}/P'$. The global proof below uses the Wannier dipole directly, so it does not require this proposition as a separate input. At $t=0$ its Wannier value is the zone average $d(0)=-(2\pi)^{-1}\int S_{\rm top}(k,0)\,dk$, in the present gauge. In general this is not $-S_{\rm top}(0,0)$.

\label{end:top}

\section{Translated Wannier frames and self-adjoint compression}\label{app:frames}
Fix $b\ge L/2+\xi$ and the slack choices of Section~\ref{sec:methods}. Set
\[
 t_{\rm fr}=\frac{e^{-\xi q/4}}{2\pi q},\quad
 \Phi^0=[-t_{\rm fr},t_{\rm fr}],\quad
 \eta'=\eta_q-\log4,\quad \eta''=\eta'-2\tau q,
 \quad a=L-4\tau.
\]
Then $2\Lambda-c(t)\le M'$ on $\Phi^0$.

\subsection{Exact translated coordinates}
For $\nu\in\Z$, let $a_\nu\in\{0,\ldots,q-1\}$ solve
$pa_\nu+\nu\equiv0\pmod q$. Put
\begin{equation}\label{eq:anchors}
 \begin{aligned}
 c_{\nu,j}&=jq+a_\nu,&
 t_{\nu,j}&=\tau_\nu+j\delta,\\
 \tau_\nu&=\theta-\nu/q+a_\nu\eps,&
 J_\nu&=\{j:|t_{\nu,j}|\le t_{\rm fr}\}.
 \end{aligned}
\end{equation}

\begin{lemma}[Covariance and well density]\label{lem:anchors}
For $c=c_{\nu,j}$ and $t=t_{\nu,j}$,
\begin{equation}\label{eq:covariance}
 S^{-c}HS^c=B_t+V_t^\eps,\quad
 V_t^\eps(m)=v(mp/q+t+m\eps)-v(mp/q+t).
\end{equation}
Each $J_\nu$ is an interval of consecutive integers. Centres in different wells have distance at least $1/(3|\delta|)$ eventually. With $y_\nu=-\tau_\nu/\delta$, the physical well bottoms are exactly
\begin{equation}\label{eq:wellpositions}
 x_\nu=qy_\nu+a_\nu=\frac{\nu-q\theta}{\delta}.
\end{equation}
Thus their density is $|\delta|$ per physical site.
\end{lemma}
\begin{proof}
Expand $(c+m)\alpha+\theta$. The terms $jp+(pa_\nu+\nu)/q$ are integers; the remaining phase is $mp/q+t+m\eps$. This proves~\eqref{eq:covariance}, including the term $a_\nu\eps$.

The affine phase $\tau_\nu+j\delta$ proves the interval assertion. If two centres belong to different wells, then
\[
 |j-j'|\,|\delta|\ge1/q-2t_{\rm fr}-2|\delta|\ge1/(2q)
\]
for large $q$. Hence $|c-c'|\ge1/(2|\delta|)-q\ge1/(3|\delta|)$. Here $q|\delta|\to0$ on the exponentially resonant scales under consideration. Substitution of $y_\nu$ gives~\eqref{eq:wellpositions}. Consecutive bottoms are separated by $1/|\delta|$, with their spatial order reversed when $\delta<0$.
\end{proof}

Translated wells and modular-inverse coordinates are classical; compare Helffer--Sj\"ostrand~\cite[Section~5.1, especially (5.1.7)--(5.1.17)]{HSII}. Here the exact coordinate identity fixes the phase shift needed for uniformity in $q$.

\subsection{A single valid local chart}
\begin{lemma}[Aubry--Perron normalization and Wannier decay]\label{lem:Wannier}
There is a real family $w_t\in\ell^2(\Z)$, $t\in\Phi^0$, such that
\begin{enumerate}
\item the translates $S^{mq}w_t$ are an orthonormal basis of the full-line top-band subspace of $B_t$;
\item $\mathcal Rw_t=w_{-t}$ and, for $0\le\ell\le3$,
\begin{equation}\label{eq:Wanniertails}
 |\partial_t^\ell w_t(n)|\le\EE e^{-a|n|};
\end{equation}
\item if $\widehat E_m(t)$ is the $m$-th Fourier coefficient of $E_t(k)$, then
\begin{equation}\label{eq:bandFourier}
 B_tw_t=\sum_m\widehat E_m(t)S^{mq}w_t,
 \qquad |\widehat E_m(t)|\le(r/2)e^{-\eta'|m|}\quad(m\ne0).
\end{equation}
These coefficients are real and even in $m$. Their first four phase derivatives are bounded by $\EE(2\pi q\Lambda)^\ell e^{-\eta'|m|}$ for $m\ne0$.
\end{enumerate}
\end{lemma}
\begin{proof}
At $t=k=0$, finite rational Aubry duality transforms $B_0(0)$ to a matrix with positive hopping $\lambda$ and diagonal $2\cos(2\pi p\ell/q)$. Adding a scalar makes it irreducible and nonnegative. Its top Perron vector has strictly positive components. The modulus of its inverse Fourier transform is maximal at original site zero, by the triangle inequality. Consequently
\[
 \gamma_0(0):=(\Pi_0(0))_{00}\ge1/q.
\]
This assertion is used only at $B_0(0)$.

Let $\gamma_t(k)=(\Pi_t(k))_{00}$, analytically continued as a matrix entry. Riesz differentiation gives
$\norm{\partial_t\Pi_t}\le4\pi\lambda r(K_q')^2$. Its $k$-Fourier coefficients are bounded by $rK_q'e^{-\eta'|m|}$. On the smaller strip $|\operatorname{Im}k|\le\eta''$, the nonconstant part is at most $4rK_q'e^{-2\tau q}$. With~\eqref{eq:slack} and $|t|\le t_{\rm fr}$, both this variation and the real-phase variation are eventually less than $1/(4q)$. Thus $\gamma_t(k)$ stays in the disc of radius $1/(2q)$ about the positive number $\gamma_0(0)$.

The principal square root is therefore analytic, and
\[
 u_t(k)=\frac{\Pi_t(k)\delta_0}{\gamma_t(k)^{1/2}}
\]
is a normalized top Bloch vector for real $k$. Uniform-gauge covariance gives $u_t(k+2\pi)=U^*u_t(k)$, $U_{ss}=e^{2\pi is/q}$, and realness gives $u_t(-k)=\overline{u_t(k)}$ on the real axis. The inverse Bloch transform
\[
 w_t(n)=\frac1{2\pi}\int_0^{2\pi}
 e^{ikn/q}u_t(k)(n\bmod q)\,dk
\]
is consequently real. The integrand is periodic in $k$. Bloch unitarity proves the orthonormal basis assertion. Shift its contour to the appropriate side of $|\operatorname{Im}k|=\eta''$ to get~\eqref{eq:Wanniertails}; $\eta''/q\ge L-4\tau$ eventually. Finitely many phase derivatives add finitely many resolvents and bounded derivatives of $D_t$, all covered by~\eqref{eq:slack}. Reflection fixes the normalization at zero and gives $\mathcal Rw_t=w_{-t}$.

Multiplication by $E_t(k)$ in the Bloch representation proves~\eqref{eq:bandFourier}. The function is even and real for real $k$. The Fourier estimate follows by contour shifting. Phase derivatives follow by differentiating the analytic top inverse in Lemma~\ref{lem:contour} and $c(t)$ up to four times; Cauchy bounds on a fixed fraction of its discriminant disc give the stated, deliberately coarse bounds.
\end{proof}

This frame is not asserted to remain anchored at the same physical site on every component of $\Phi_{\rm fr}$ modulo $1/q$. Lemma~\ref{lem:anchors} translates each well before the local construction is applied.

\subsection{The global Gram matrix}
Let $\cF=\{(\nu,j):j\in J_\nu\}$ and define
\[
 \widehat w_f=S^{c_f}w_{t_f},\quad
 W\beta=\sum_{f\in\cF}\beta_f\widehat w_f,
 \quad \Gamma_{f'f}=\inner{\widehat w_{f'}}{\widehat w_f}.
\]

\begin{lemma}[Overlap cancellation and global orthogonalization]\label{lem:Gram}
Put
\[
 g_\infty=\exp\left(-\frac{L-6\tau}{3|\delta|}\right),\qquad
 g_1=|\delta|e^{-(L-5\tau)q}+g_\infty.
\]
Then $\norm{\Gamma-1}\le g_1\le1/4$ eventually. The closed span $\cV=\overline{\Ran W}$ has orthogonal projection
\[
 \Pi=W\Gamma^{-1}W^*,\qquad
 U_V=W\Gamma^{-1/2}:\ell^2(\cF)\longrightarrow\cV
\]
unitary, and $\norm{\Gamma^{-1/2}-1}\le g_1$, $\norm{\Gamma^{-1}}\le4/3$.
\end{lemma}
\begin{proof}
For equal phase, $\inner{S^{nq}w_t}{w_t}=\delta_{n0}$. For $n\ne0$, use a mean-value integral in the phase before summing spatially:
\[
 |\inner{S^{nq}w_{t'}}{w_t}|
 \le\EE|t'-t|(1+|n|q)e^{-a|n|q}.
\]
For $n=0$, normalized real vectors give
$\inner{w_{t'}}{w_t}-1=-\norm{w_{t'}-w_t}^2/2$.
The elementary convolution estimate used here is
\[
 \sum_x e^{-a|x-d|}e^{-a|x|}
 \le(|d|+1+2/(1-e^{-2a}))e^{-a|d|}.
\]
In one well $t'-t=(j'-j)\delta$. Its off-diagonal Schur sums are at most $|\delta|e^{-(L-5\tau)q}$. Between wells the physical separation is at least $1/(3|\delta|)$; summing over the distinct integer centres gives at most $g_\infty$. The row and column constants are independent of the number of wells. This proves the norm bound.

Spectral calculus on $[3/4,5/4]$ gives the inverse and square-root bounds. The identities follow from $W^*W=\Gamma$. This is the usual Gram orthogonalization used in multiple-well analysis~\cite{HSwells}.
\end{proof}

\subsection{The matrix and its quantization convention}
The local Taylor bound is
\begin{equation}\label{eq:Taylor}
 |v(x+y)-v(x)-yv'(x)|\le4\pi^2\lambda y^2.
\end{equation}
Let $M=W^*HW$. In a well define the column-frozen matrix
\[
 h_{\rm tb}(j',j)=\widehat E_{j'-j}(t_j)
                  +\delta_{jj'}s(t_j),
 \qquad s(t)=\eps\inner{w_t}{XD_tw_t}.
\]
Set its entries between different wells to zero and put $h_s=(h_{\rm tb}+h_{\rm tb}^*)/2$.

\begin{proposition}[Self-adjoint band compression]\label{prop:compression}
The function $d(t)=s(t)/\eps$ is real analytic and even on $\Phi^0$, with $|\partial_t^\ell d|\le\EE$ for $\ell\le3$. Moreover
\begin{align}
 \norm{M-h_{\rm tb}}&\le\eta_1
  :=|\delta|e^{-(L-6\tau)q}+\delta^2\EE+Cg_\infty,\label{eq:eta1}\\
 \norm{\widetilde M-h_s}&\le\eta_2
  \le |\delta|e^{-(L-7\tau)q}+\delta^2\EE,
 \quad \widetilde M=\Gamma^{-1/2}M\Gamma^{-1/2}.
 \label{eq:eta2}
\end{align}
The latter is the matrix of $\Pi H\Pi|_{\cV}$ in an orthonormal basis.
\end{proposition}
\begin{proof}
Use~\eqref{eq:covariance} and expand $V_t^\eps=\eps XD_t+\rho_t$ with $|\rho_t(m)|\le4\pi^2\lambda m^2\eps^2$. In one well let $n=j'-j$, $t'=t+n\delta$. The frozen part of $M_{j'j}$ is
\[
 \sum_m\widehat E_m(t)
 \inner{S^{nq}w_{t'}}{S^{mq}w_t}.
\]
At $n=0$ this is exactly $\widehat E_0(t)$. At $n\ne0$, subtract $\widehat E_n(t)$, use the overlap bound in Lemma~\ref{lem:Gram}, and convolve the two exponential tails. Its Schur sum is at most $|\delta|e^{-(L-6\tau)q}$ after absorbing polynomial factors.

The diagonal first-order term is exactly $s(t)$. For $n\ne0$, its absolute entry is at most
$|\eps|\EE(1+|n|q)^2e^{-a|n|q}$. The Taylor remainder has entries bounded by
$\eps^2\EE(1+|n|q)^3e^{-a|n|q}$, including the diagonal. Their Schur sums give the two terms in~\eqref{eq:eta1}. Cross-well entries are at most a constant times the overlap tail, hence cost $Cg_\infty$.

Reflection changes the signs of $X$ and $D_t$ simultaneously: $\mathcal RX\mathcal R=-X$ and $\mathcal RD_t\mathcal R=-D_{-t}$. Lemma~\ref{lem:Wannier} implies $d(-t)=d(t)$; its derivatives are bounded by the summable site tails.

For orthogonalization use $\norm M\le\norm\Gamma\norm H$, valid for a nonorthogonal frame. Thus
\[
 \norm{\widetilde M-M}\le3g_1(3+2\lambda).
\]
Taking the Hermitian part of the bound against $h_{\rm tb}$ proves~\eqref{eq:eta2}. Every subsequent min--max statement uses $h_s$, which is self-adjoint.
\end{proof}

For $a(t,k)=\sum_n\widehat a_n(t)e^{ink}$ with real even coefficients, define the symmetric quantization on the lattice $t_j=\tau+ j\delta$ by
\begin{equation}\label{eq:quantization}
 \Op^\tau(a)(j',j)
 =\tfrac12\bigl(\widehat a_{j'-j}(t_j)
                +\widehat a_{j'-j}(t_{j'})\bigr).
\end{equation}
Then $h_s=\bigoplus_\nu\Op^{\tau_\nu}(E_t(k)+s(t))|_{J_\nu}$. In particular
\[
 \Op^\tau(g(k))=g(S),\qquad
 \Op^\tau(u(t)g(k))=\tfrac12(ug(S)+g(S)u).
\]
This mean of row- and column-frozen quantizations is the declared convention throughout. No Weyl-quantization identity is silently substituted for it.

\label{end:frames}

\section{Global exclusion, boundary weights and exact box counts}\label{app:global}
Write $Q=1-\Pi$. The physical translates of every well now belong to one global frame.

\subsection{Excluding the complementary subspace}
Set
\begin{equation}\label{eq:m0}
 M''=\frac4{9\pi^2}\Lambda e^{-\xi q/2},\qquad
 m_0=\frac{M''}{2P_*'}
 \ge\frac{e^{-(\xi/2+h-L)q}}{Cq}.
\end{equation}

\begin{lemma}[Complement exclusion]\label{lem:complement}
For all sufficiently large $q$,
\begin{equation}\label{eq:complement}
 QHQ\le(E_*-m_0)Q.
\end{equation}
The same form inequality holds for vectors supported in a box with ends midway between well bottoms and orthogonal to the frames of its internal wells.
\end{lemma}
\begin{proof}
Put $\ell_c=\lceil(q|\delta|)^{-1/2}\rceil$, $\ell=\ell_cq$. Choose a squared partition $\sum_a\chi_a^2=1$ with support intervals $I_a$ of at most $4\ell$ sites, overlap at most three, and Lipschitz constants at most $\pi/\ell$. Discrete IMS has error at most $3\pi^2/\ell^2\le3\pi^2|\delta|/q$. Freezing the phase on a patch costs at most
\[
 \mathrm{drift}_a\le16\pi\lambda\sqrt{|\delta|/q}.
\]
Let $j_a$ be its centre cell.

If some translated coordinate satisfies $|t_{\nu,j_a}|\le t_{\rm fr}/2$, all cells meeting the patch belong to $J_\nu$ for large $q$. Indeed their phase drift is at most $3\ell_c|\delta|$, which is much less than $t_{\rm fr}$. Freeze the reference $B_{t_a}$ at $c_a=c_{\nu,j_a}$ and write $f_a=\chi_a\psi$. Its other bands lie below $E_*-|I_q|$, so
\begin{equation}\label{eq:localcomplement}
 \inner{f_a}{Hf_a}\le(E_*-|I_q|+\mathrm{drift}_a)\norm{f_a}^2
                   +|I_q|\norm{\Pi_af_a}^2,
\end{equation}
where $\Pi_a$ is its frozen top-band projection.

We give the joint Bessel estimate, including its support factors. For $z=c_a+mq$ define
\[
 g_{a,m}=(\chi_a-\chi_a(z))S^zw_{t_a}
       +\chi_a(z)S^z(w_{t_a}-w_{t_a+m\delta}).
\]
The second term is defined only when its coefficient is nonzero; then both phases lie in the chart. Orthogonality to the translated frame implies
$\inner{S^zw_{t_a}}{\chi_a\psi}=\inner{g_{a,m}}{\psi}$. With $\zeta_q=\EE\sqrt{|\delta|/q}$, the tails and the mean-value phase bound give
\begin{equation}\label{eq:supportBessel}
 |g_{a,m}(x)|\le C\zeta_q(1+|x-z|)e^{-a|x-z|}
       \bigl(\un_{I_a}(x)+\un_{I_a}(z)\bigr).
\end{equation}
For fixed $x$, the first support term has at most three patches. In the second, each centre belongs to its unique well and to at most three patch supports. Summation over the distinct integer centres is bounded by the summable exponential envelope. Both Schur sums of the synthesis kernel are at most $C\zeta_q$, including the column sum over $x$. Therefore
\begin{equation}\label{eq:jointBessel}
 \sum_{a\text{ of this type}}\norm{\Pi_a\chi_a\psi}^2
 \le\EE\frac{|\delta|}{q}\norm\psi^2.
\end{equation}
The support factors in~\eqref{eq:supportBessel} are essential; an unrestricted sum of envelopes over all patches would not justify~\eqref{eq:jointBessel}.

On the other patches, the frozen phase has distance at least $t_{\rm fr}/3$ from $q^{-1}\Z$ throughout the patch. Thus $2\Lambda-c\ge M''$. Convexity on the top branch yields the sharp drop $M''/P_*'$ from $E_*$. Below that branch the drop is at least $|I_q|\ge2R/P_*'$. Hence
\[
 \inner{f_a}{Hf_a}\le(E_*-2m_0+\mathrm{drift}_a)\norm{f_a}^2.
\]
Summing the two patch types gives an error bounded by
\[
 C\sqrt{|\delta|/q}+C|\delta|/q
       +|I_q|\EE|\delta|/q.
\]
It is $o(m_0)$ uniformly on $b\ge L/2+\xi$, after~\eqref{eq:slack}. Also $2m_0\le|I_q|$. This proves~\eqref{eq:complement}. There is no additive $e^{-q}$ floor.

For a supported box vector, every patch of the first type that meets the box is associated with an internal well. The patch length is $o(1/|\delta|)$ and the ends lie in the middle of the forbidden region. The same orthogonality and sums therefore apply using only the internal frames.
\end{proof}

\subsection{Coupling and the actual boundary weight}
For $f=(\nu,j)$ let $d_f=\dist(j,\Z\setminus J_\nu)\ge1$ and
$\Delta_{\rm bd}=\operatorname{diag}(e^{-\eta'd_f})$.

\begin{lemma}[Interior coupling and omitted translates]\label{lem:coupling}
There is a decomposition
\[
 QHW\beta=Q\sum_f\beta_f(z_f^{\rm in}+z_f^{\rm bd})
\]
with
\begin{align}
 \left\|\sum_f\beta_fz_f^{\rm in}\right\|&\le c_{\rm in}\norm\beta,
  &c_{\rm in}&=|\delta|\EE,\label{eq:interiorcoupling}\\
 \left\|\sum_f\beta_fz_f^{\rm bd}\right\|^2
 &\le16r^2e^{-\eta'}\inner\beta{\Delta_{\rm bd}\beta}
       +Cr^2g_\infty\norm\beta^2.\label{eq:boundaryform}
\end{align}
\end{lemma}
\begin{proof}
Use~\eqref{eq:bandFourier} in~\eqref{eq:covariance}. For $j+m\in J_\nu$ subtract the retained frame vector, which is killed by $Q$. This leaves
\[
\begin{split}
 z_f^{\rm in}={}&\sum_{m\ne0,\,j+m\in J_\nu}\widehat E_m(t_f)
       S^{c_f+mq}(w_{t_f}-w_{t_f+m\delta})
       +S^{c_f}(\eps XD_{t_f}+\rho_{t_f})w_{t_f},\\
 z_f^{\rm bd}={}&\sum_{j+m\notin J_\nu}\widehat E_m(t_f)
                         S^{c_f+mq}w_{t_f}.
\end{split}
\]
The first family has envelope
$|\delta|\EE(1+|x-c_f|)^2e^{-a|x-c_f|}$.
Schur summation over distinct centres gives~\eqref{eq:interiorcoupling}. In the second family only $|m|\ge d_f$ occurs, so
\[
 \norm{z_f^{\rm bd}}\le2re^{-\eta'd_f},\qquad
 |z_f^{\rm bd}(x)|\le r\EE e^{-2\tau qd_f}e^{-a|x-c_f|}.
\]
Within one well use weighted Cauchy--Schwarz and
$\sum_{f\text{ in well}}\norm{z_f^{\rm bd}}\le8re^{-\eta'}$.
This gives the first term in~\eqref{eq:boundaryform}. The convolution of the physical-site envelopes between distinct wells has Schur sum at most $Cr^2g_\infty$, which gives the second.
\end{proof}

Define
\begin{equation}\label{eq:sandwichdefs}
 \begin{split}
 h^-&=h_s-\eta_2,\qquad h^+=h_s+\eta_3+W_{\rm bd},\\
 \eta_3&=\eta_2+\frac{16c_{\rm in}^2}{3m_0}
                 +\frac{Cr^2g_\infty}{m_0},\\
 W_{\rm bd}&=\frac{64r^2e^{-\eta'}}{m_0}
           \Gamma^{-1/2}\Delta_{\rm bd}\Gamma^{-1/2}\ge0.
 \end{split}
\end{equation}

\begin{lemma}[Two-sided form comparison]\label{lem:sandwich}
On $\cV$ the compression is at least $U_Vh^-U_V^*$. On the whole physical space,
\begin{equation}\label{eq:sandwich}
 H\le U_Vh^+U_V^*\Pi+(E_*-m_0/2)Q.
\end{equation}
In rescaled units the errors satisfy
\begin{align}
 \frac{P_*'\eta_2}{\hbar}
 &\le e^{-(L/2-8\tau)q}+e^{-(b-L/2)q}\EE,\label{eq:errors1}\\
 \frac{P_*'c_{\rm in}^2}{m_0\hbar}
 &\le e^{-(b-L/2-\xi/2)q}\EE,\label{eq:errors2}\\
 P_*'\norm{W_{\rm bd}}
 &\le e^{-(L-\xi/2-3\tau)q}.\label{eq:errors3}
\end{align}
\end{lemma}
\begin{proof}
The lower assertion is Proposition~\ref{prop:compression}. Write a vector as $U_V\beta+y$, $y\in\Ran Q$, and apply
\[
 2\operatorname{Re}\inner y{QHU_V\beta}
 \le\frac{m_0}{2}\norm y^2+\frac2{m_0}\norm{QHU_V\beta}^2.
\]
Insert Lemmas~\ref{lem:complement} and~\ref{lem:coupling}, with coefficients $\Gamma^{-1/2}\beta$, whose squared norm is at most $4\norm\beta^2/3$. The constants in~\eqref{eq:sandwichdefs} dominate the resulting terms.

For the sizes use $P_*'\le C_hqe^{hq}$, $r\le(\lambda+1)/q$, $P_*'r\le C\Lambda\EE$, and
$e^{-\eta'}=32q\Lambda^{-1}e^{\tau q}$. Equations~\eqref{eq:eta2} and~\eqref{eq:m0} give~\eqref{eq:errors1}--\eqref{eq:errors2}; the same substitutions give~\eqref{eq:errors3}, after the prescribed slack. All omitted interwell tails are smaller uniformly in $b$.
\end{proof}

The whole norm $\norm{W_{\rm bd}}/\hbar$ need not tend to zero when $b>3L/2$. Its distance from the low-energy part of a well must be retained.

\begin{lemma}[Boundary leakage with the correct rate]\label{lem:leakage}
Let $K(d_*)=\{f:d_f\ge d_*\}$. Then
\begin{equation}\label{eq:weightedleakage}
 P_*'\norm{\un_KW_{\rm bd}\un_K}
 \le e^{(\xi/2+3\tau)q}
       \bigl(Ce^{-\eta'd_*/2}+Cg_\infty\bigr)^2.
\end{equation}
For the inner region $|2\pi q\sqrt\Lambda\,t_f|\le4$, its right-hand side is smaller than every fixed power of $\hbar$, uniformly in $b\ge L/2+\xi$.
\end{lemma}
\begin{proof}
Let $\Gamma_0$ be the direct sum of the well blocks. In each block use the weighted Schur algebra with weight $e^{\eta'|j-j'|/2}$. Lemma~\ref{lem:Gram} gives $\norm{\Gamma_0-1}_{\rm weighted}=o(1)$. The convergent binomial series gives the same bounded weighted norm for $\Gamma_0^{-1/2}$. Since $d_j$ is 1-Lipschitz,
\[
 e^{-\eta'd_g/2}\le e^{-\eta'd_f/2}e^{\eta'|g-f|/2}.
\]
Both restricted Schur sums show
$\norm{\Delta_{\rm bd}^{1/2}\Gamma_0^{-1/2}\un_K}
\le Ce^{-\eta'd_*/2}$. Analytic functional calculus and the interwell Gram bound give
$\norm{\Gamma^{-1/2}-\Gamma_0^{-1/2}}\le Cg_\infty$.
Finally $P_*'64r^2e^{-\eta'}/m_0\le e^{(\xi/2+3\tau)q}$, proving~\eqref{eq:weightedleakage}.

The inner region lies at distance
\[
 d_*\ge\frac{e^{(b-\xi/4)q}}{8\pi q}
\]
from the frame boundary. Since $\eta'$ is of order $q$, its leakage has rate
\begin{equation}\label{eq:doubleexp}
 \exp\bigl(-c e^{(b-\xi/4)q}\bigr),
\end{equation}
with the still smaller interwell term retained in~\eqref{eq:weightedleakage}. There is no extra factor $q$ in the exponent of~\eqref{eq:doubleexp}. For any fixed $N$, $e^{(b-\xi/4)q}$ dominates $Nbq$ uniformly for $b\ge L/2+\xi$. As $\hbar=2\pi qe^{-(b-L/2)q}$, the asserted comparison follows.
\end{proof}

\subsection{Arbitrary intervals and whole-well boxes}
\begin{inputbox}
\begin{classicalinput}[IDS limit; Avron--Simon]
For fixed irrational $\alpha$ and fixed phase, every sequence of intervals $I_N\subset\Z$ with $|I_N|\to\infty$ satisfies
\[
 \frac{\#\{\operatorname{eig}H_{I_N}\le E\}}{|I_N|}\longrightarrow N(E)
 \quad(E\in\R),
\]
independently of the positions of the intervals.
\end{classicalinput}
\end{inputbox}
\begin{proof}
For each fixed $k$, finite propagation bounds the difference between $\tr H_I^k$ and $\sum_{j\in I}\inner{\delta_j}{H^k\delta_j}$ by a constant depending on $k$ and $\norm H$, independent of $I$. The latter diagonal is a continuous trigonometric polynomial of $\theta+j\alpha$. Unique ergodicity gives its interval average uniformly in the starting point. Thus all moments of the normalized counting measures converge to the DOS moments. Uniform compact spectral support and polynomial approximation imply weak convergence. Continuity of the irrational IDS gives convergence of the distribution functions at every $E$. This recovers the classical IDS passage~\cite{AvronSimon} in precisely the moving-box form needed here.
\end{proof}

Choose box ends at the integer parts of the midpoints between the $x_\nu$ in~\eqref{eq:wellpositions}, with spatial orientation reversed if $\delta<0$. A box $I$ containing $K$ consecutive wells has
\begin{equation}\label{eq:boxlength}
 |I|=K/|\delta|+O(1).
\end{equation}
Every frame centre in it is at distance at least $1/(3|\delta|)$ from its exterior. Choose $c_{\rm box}>0$ slightly smaller than $a/3$ and set
$g_{\rm box}=\exp(-c_{\rm box}/|\delta|)$, decreasing $c_{\rm box}$ further when needed to absorb polynomial and $m_0^{-1}$ factors.

For each well let $\Gamma_\nu$ be its Gram block and define
\begin{equation}\label{eq:boxoperators}
 \begin{split}
 h_\nu^-&=h_{s,\nu}-\eta_2',\qquad
 h_\nu^+=h_{s,\nu}+\eta_3'+W_{{\rm bd},\nu},\\
 W_{{\rm bd},\nu}&=\frac{64r^2e^{-\eta'}}{m_0}
                 \Gamma_\nu^{-1/2}\Delta_\nu\Gamma_\nu^{-1/2},\\
 \eta_2'&=\eta_2+Cg_{\rm box},\\
 \eta_3'&=\eta_3+Cg_{\rm box}
       +C(P_*')^{-1}e^{(\xi/2+3\tau)q}g_\infty.
 \end{split}
\end{equation}
Constants may be enlarged once, independently of $K$.

\begin{proposition}[Exact box sandwich]\label{prop:box}
For $E>E_*-m_0/2$,
\begin{equation}\label{eq:boxsandwich}
 \sum_{\nu\text{ in }I}\#\{\operatorname{eig}h_\nu^->E\}
 \le\#\{\operatorname{eig}H_I>E\}
 \le\sum_{\nu\text{ in }I}\#\{\operatorname{eig}h_\nu^+>E\}.
\end{equation}
If both per-well counts equal $D$ at every well, then the middle count is exactly $KD$, and $N(E)=1-D|\delta|$.
\end{proposition}
\begin{proof}
Truncate the internal frame vectors by $\un_I$. Their Gram and compression matrices differ from the internal blocks of $\Gamma$ and $M$ in Schur norm by at most $Cg_{\rm box}$. The tail entries contain the distances from both centres to a box end. Summing them over the distinct physical centres is independent of $K$; there are only two ends.

Orthogonalize the truncated frame. A box vector orthogonal to it is orthogonal to the original internal frames, because it vanishes outside $I$. Lemma~\ref{lem:complement} applies. In the coupling calculation all retained frozen translates still belong to internal wells and are killed by the box complementary projection. The omitted truncation term costs $Cg_{\rm box}$ before squaring. Retain its square divided by $m_0$ and its mixed terms until absorption: $g_{\rm box}/m_0=o(1)$ and $(c_{\rm in}+re^{-\eta'})/m_0=o(1)$, so an enlarged $Cg_{\rm box}$ covers them. No constant grows with $K$.

Separating the well Gram blocks from the full Gram matrix costs $Cg_\infty$. Conjugation of the boundary weight costs at most the last term of $\eta_3'$ in~\eqref{eq:boxoperators}. Thus the form comparison of Lemma~\ref{lem:sandwich} holds in the box with these already enlarged operators. Interlacing for its compression gives the lower count. The form majorant is block diagonal with complementary block at most $E_*-m_0/2<E$, giving the upper count by min--max. This proves~\eqref{eq:boxsandwich}.

If both sums equal $KD$, the integer inequalities force equality exactly. Let $K\to\infty$ at the fixed scale. Equation~\eqref{eq:boxlength} and the classical IDS input above yield $1-N(E)=D|\delta|$.
\end{proof}

The final gap inference uses full DOS support. For completeness, if a nonnegative continuous function supported in an interval has DOS integral zero, then the phase average of its nonnegative diagonal matrix element is zero. That element is continuous in phase, hence vanishes at every phase. Covariance gives the same assertion at every site. The positive operator is zero, so its support misses the common spectrum. This proves the required support statement for the minimal continuous AMO family.

\label{end:global}

\section{Common well windows and completion of the scale theorem}\label{app:comparison}
Let $T(k)=2-2\cos k$, $T(S)=2-S-S^*$, and
\[
 U(t)=2\Lambda(1-\cos2\pi qt),\qquad
 F(y)=P_*'\bigl(E_*-P_{\rm top}^{-1}(R-y)\bigr).
\]
The inverse is the one selected in Lemma~\ref{lem:contour}. On each well use
\[
 z_j=2\pi q\sqrt\Lambda\,t_{\nu,j}
     =\hbar\sgn\delta\,(j-y_\nu),\qquad
 I_{\rm in}=\{j:|z_j|\le3.5\},\quad I_{\rm out}=\{j:|z_j|\ge3\}.
\]
The frame extends to $|z|=\sqrt\Lambda e^{-\xi q/4}\to\infty$.

Define the rescaled self-adjoint operators
\begin{equation}\label{eq:Aoperators}
 A_\nu=P_*'(E_*+s(0)-h_{s,\nu}),\qquad
 A_\nu^{(\pm)}=P_*'(E_*+s(0)-h_\nu^\pm).
\end{equation}
The superscript follows $h_\nu^\pm$, so $A_\nu^{(+)}$ is the lower rescaled comparison. Chambers gives exactly
\begin{equation}\label{eq:Asymbol}
 A_\nu=\Op^{\tau_\nu}(F(T(k)+U(t)))
                    -P_*'(s(t)-s(0)).
\end{equation}
This is the actual top inverse branch; no branch is recovered from a collapsed polynomial spectrum.

\subsection{Inverse curvature and a two-region count}
Cauchy bounds on a fixed fraction of $|y|\le R/(8q)$ give
\begin{equation}\label{eq:fcurvature}
 f:=\sup_{|y|\le2M'}|F''(y)|\le\EE/\Lambda.
\end{equation}
Indeed $|F|\le P_*'r/2$ there, and $P_*'r\le\Lambda\EE$, for example by Markov's inequality on $I_q$. We have $F(0)=0$, $F'(0)=1$. For the nonnegative real arguments used below, convexity gives $F(y)\ge y$ and $1\le F'(y)\le1+2fy$ within the relevant interval. On a complex disc only the bound $|F'(z)-1|\le f|z|$ is used; no ordered complex inequality is intended.

\begin{lemma}[Two-region min--max]\label{lem:tworegion}
Suppose a self-adjoint matrix $A$ on an interval containing $I_{\rm in}$ has
\[
 \sup_j\sum_n|A(j,j+n)|n^2\le C_A
\]
and its restriction to $I_{\rm out}$ has form lower bound six. For $e\le5$ and sufficiently small $\kappa$,
\begin{equation}\label{eq:tworegion}
 \#\{\operatorname{eig}A_{I_{\rm in}}<e\}
 \le\#\{\operatorname{eig}A<e\}
 \le\#\{\operatorname{eig}A_{I_{\rm in}}<e+16C_A\kappa\}.
\end{equation}
\end{lemma}
\begin{proof}
Choose $\chi_{\rm in}^2+\chi_{\rm out}^2=1$, with $\chi_{\rm in}=1$ on $|z|\le3$, supported in $|z|\le3.5$, and Lipschitz constants at most four in $z$. Their lattice Lipschitz constants are at most $4\hbar$. The IMS double-commutator Schur norm is at most $16C_A\kappa$. The map $J\psi=(\chi_{\rm in}\psi,\chi_{\rm out}\psi)$ is an isometry into the doubled space, and
\[
 A\ge J^*(A_{I_{\rm in}}\oplus6I)J-16C_A\kappa.
\]
Compression to the image of this isometry can only raise the ordered lower eigenvalues. Below $e+16C_A\kappa<6$, the doubled-space count comes entirely from the first block. This proves the upper inequality. The lower one is ordinary interlacing for the inner compression.
\end{proof}

\begin{lemma}[Far barrier and row moments]\label{lem:far}
Both $A_\nu^{(\pm)}$ satisfy Lemma~\ref{lem:tworegion} with $C_A=4$, uniformly in every well and $b\ge L/2+\xi$.
\end{lemma}
\begin{proof}
For $|\operatorname{Im}k|\le1$, expand
\[
 F(U+2-2\cos k)=F(U+2)-2\cos k\,F'(U+2)+O(f).
\]
Its nearest off-diagonal coefficients have absolute value at most $1+2fM'+2f$, and the remaining coefficients are at most $15fe^{-|n|}$. As $fM'\to0$, the off-diagonal row sums and second moments are at most $2.01+o(1)$, even after the symmetric row/column average.

The diagonal average of $F(T+U)$ is at least $2+U$. Evenness of $s$ and~\eqref{eq:dipole} give
\begin{equation}\label{eq:dipolerelative}
 |P_*'(s(t)-s(0))|\le|\delta|\EE U(t).
\end{equation}
On $I_{\rm out}$, $U\ge8.9$ eventually. Subtracting the row sum leaves a form barrier at least $8.7$. The scalar comparison errors in~\eqref{eq:boxoperators} are $o(1)$ in rescaled units. So is the whole norm of $P_*'W_{{\rm bd},\nu}$ by~\eqref{eq:errors3}. Its second row moments are $o(1)$ as well, by the weighted Gram series. The barrier remains above six, and the moments below four.
\end{proof}

The Mathieu operator $\hM_{\kappa,y}=T(S)+\kappa(j-y)^2$ satisfies the same lemma with $C_A=2$.

\subsection{Relative comparison inside the well}
\begin{proposition}[The two relative errors]\label{prop:relative}
On $I_{\rm in}$, for either sign,
\begin{equation}\label{eq:relativeform}
 (1-\eta_r)\hM_{\kappa,y_\nu}|_{I_{\rm in}}-\eta_a
 \le (A_\nu^{(\pm)})_{I_{\rm in}}
 \le(1+\eta_r)\hM_{\kappa,y_\nu}|_{I_{\rm in}}+\eta_a,
\end{equation}
where
\begin{align}
 \eta_r&\le C(f+\Lambda^{-1}+|\delta|\EE),\label{eq:etar}\\
 \eta_a&\le C\bigl(f\kappa+P_*'\eta_2'+P_*'\eta_3'
       +P_*'\norm{\un_{I_{\rm in}}W_{{\rm bd},\nu}\un_{I_{\rm in}}}\bigr).
 \label{eq:etaa}
\end{align}
For some $c>0$, $\eta_a/\hbar\le e^{-cq}$ uniformly in $b$. One may choose any sufficiently small positive $c$ below $\min\{\xi/2,L/2-8\tau\}$ after the rate losses.
\end{proposition}
\begin{proof}
For the symbol write
\[
 F(T+U)=F(T)+UF'(T)+U^2\varrho(t,k),\qquad
 \varrho=\int_0^1(1-r)F''(T+rU)\,dr.
\]
On the complex unit strip in $k$ over $I_{\rm in}$, $|\varrho|\le Cf$. First take pure multiplier inequalities on the full line, then compress. Since $0\le T(S)\le4$,
\[
 T(S)\le F(T(S))\le(1+Cf)T(S).
\]
For $g=F'(T(S))$ the declared symmetric quantization gives
\[
 \tfrac12(Ug+gU)=U^{1/2}gU^{1/2}
                 +\tfrac12[U^{1/2},[U^{1/2},g-1]].
\]
The first term is between $U$ and $(1+Cf)U$. The function $U^{1/2}$ is $\hbar$-Lipschitz per cell. The Fourier coefficients of $g-1$ have second moment $O(f)$, so the double commutator has norm at most $Cf\kappa$.

The coefficients of $\varrho$ are bounded by $Cfe^{-|n|}$. On the inner region $U\le12.3$ and
$U_{j+n}\le2U_j+2\kappa n^2$. Schur's quadratic-form estimate then bounds the quantized $U^2\varrho$ by $CfU+Cf\kappa$ in absolute form. The cosine Taylor bound gives
\[
 (1-1.1/\Lambda)z_j^2\le U(t_j)\le z_j^2
          =\kappa(j-y_\nu)^2.
\]
The dipole contributes the relative term in~\eqref{eq:dipolerelative}. Adding the scalar errors and the compressed boundary weight proves~\eqref{eq:relativeform}--\eqref{eq:etaa}.

For the sizes, $f\kappa/\hbar=f\hbar\to0$ exponentially. Equations~\eqref{eq:errors1}--\eqref{eq:errors2} control the scalar errors divided by $\hbar$. The box and interwell tails are smaller than every power of $\hbar$, as is the inner boundary leakage by Lemma~\ref{lem:leakage}. This proves the stated uniform rate. We have used compression of $F(T_{\Z})$; it is not identified with $F(T_{I_{\rm in}})$.
\end{proof}

The term $|\delta|\EE$ in~\eqref{eq:etar} need not be $e^{-(L-\xi/2)q}$. At $b=L/2+\xi$ with small $\xi$ such an upper bound would be false. Its contribution is controlled with the other cap.

\begin{proposition}[One count in every offset]\label{prop:commoncount}
Let $c_1$ be sufficiently small for the Mathieu ladder~\eqref{eq:Mladder} of Appendix~\ref{app:mathieu}. For every integer $D$ satisfying~\eqref{eq:Dcap}, every well $\nu$ and every $e\in W_D$ from~\eqref{eq:WD},
\begin{equation}\label{eq:commoncount}
 \#\{\operatorname{eig}A_\nu^{(+)}<e\}
 =\#\{\operatorname{eig}A_\nu^{(-)}<e\}=D.
\end{equation}
The threshold in $q$ is independent of $b,D,\nu$ and $\theta$.
\end{proposition}
\begin{proof}
Apply Lemma~\ref{lem:tworegion} first to $A=A_\nu^{(\pm)}$, then use~\eqref{eq:relativeform}, and then apply it to the full-line Mathieu operator. The count is squeezed between its counts at
\[
 x_\pm=e\pm(2\eta_re+2\eta_a+128\kappa).
\]
Here $e\le C_M(D+1)\hbar$ and the cap makes $e$ stay below the fixed cutoff of the Mathieu input in Appendix~\ref{app:mathieu}. Crucially,
\begin{align}
 D(f+\Lambda^{-1})&\le e^{-(\xi-o(1))q},\label{eq:cap1use}\\
 D|\delta|\EE&\le\frac{c_1|\delta|\EE}{\hbar}
        =\frac{c_1}{2\pi q}e^{-(L/2-o(1))q}.
 \label{eq:cap2use}
\end{align}
The second cap on $D$ proves~\eqref{eq:cap1use}; the first proves~\eqref{eq:cap2use}. Hence $(D+1)\eta_r=o(1)$ uniformly. Also $\eta_a/\hbar=o(1)$ and $\kappa/\hbar=\hbar=o(1)$. The points $x_\pm$ therefore stay inside
\[
 [\mu_{D-1}+c_M\hbar/8,\ \mu_D-c_M\hbar/8]
\]
with strict interior margins for large $q$. The Mathieu input~\eqref{eq:Mladder} gives exactly $D$ levels below each, whatever the offset $y_\nu$. This proves~\eqref{eq:commoncount}. All errors were uniform over the declared parameters.
\end{proof}

\begin{proof}[Completion of Theorem~\ref{thm:scale}]
For $E$ corresponding to the interior of $W_D$, its rescaled energy is bounded and $|s(0)|\le|\delta|\EE$. Equations~\eqref{eq:topbounds} and~\eqref{eq:m0} show $E>E_*-m_0/2$ eventually. Proposition~\ref{prop:commoncount} applies to the already enlarged box operators~\eqref{eq:boxoperators}. Proposition~\ref{prop:box} therefore gives~\eqref{eq:IDSidentity}. The argument in Section~\ref{sec:consequences} proves the label and the width with $c_0=c_M/2$, or a smaller absolute constant.

Finally
\[
 \frac{c_0\hbar}{P_*'}\ge C_he^{-(b+L/2+(h-L))q}.
\]
The finite derivative list permits $h-L$ as small as desired. Diagonalize the corresponding thresholds for the $o(1)$ formulation; every bound is uniform in $b$. If $b\le3L/2-\xi$, then
$c_1/\hbar=c_1e^{(b-L/2)q}/(2\pi q)\le e^{(L-\xi)q}$ eventually, proving the cap assertion. No parity condition entered the construction.
\end{proof}

\label{end:comparison}

\section{The uniform Mathieu ladder: a one-well argument}\label{app:mathieu}
This is input (I5), stated in the form used by the comparison. The semiclassical facts are classical. We give the one-well reduction, the oscillator endpoint and the twist bookkeeping, since a fixed-index expansion alone does not control a growing ladder.

\begin{inputbox}
\begin{classicalinput}[uniform scalar Mathieu ladder; one-well quantization]
There are absolute $c_M,C_M,c_1,\hbar_0>0$ such that, for $0<\hbar<\hbar_0$, $\kappa=\hbar^2$, the simple eigenvalues $e_j(\kappa,y)$ of
\[
 \hM_{\kappa,y}=2-S-S^*+\kappa(j-y)^2
 \quad\hbox{on }\ell^2(\Z)
\]
have $y$-independent reference levels $\mu_j(\kappa)$ with
\begin{equation}\label{eq:Mladder}
 \begin{split}
 |e_j(\kappa,y)-\mu_j(\kappa)|&\le c_M\hbar/8,\\
 \mu_{j+1}-\mu_j&\ge c_M\hbar,\qquad
 \mu_j\le C_M(j+1)\hbar.
 \end{split}
\end{equation}
These estimates are uniform for all $y\in\R$ and all levels in a fixed compact energy interval below four, in particular for all $j\le c_1/\hbar$ with $\mu_j\le3.5$ after decreasing $c_1$. The same common windows separate the union over offsets.
\end{classicalinput}
\end{inputbox}
\begin{proof}
\emph{Fourier transformation.} Use $\widehat u(k)=\sum_ju_je^{ijk}$. Position becomes $-i\partial_k$ and hopping becomes $2-2\cos k$. Writing $\widehat u=e^{iyk}g$ gives
\begin{equation}\label{eq:circleMathieu}
 -\hbar^2g''+(2-2\cos k)g,
 \quad g(\pi)=e^{-2\pi iy}g(-\pi),\quad
 g'(\pi)=e^{-2\pi iy}g'(-\pi).
\end{equation}
Thus the offset is solely a circle twist. Integer changes of $y$ do not affect the spectrum.

\emph{A well separated from the seam.} Fix energies $3.5<e_+<4$. The allowed set for $V(k)=2-2\cos k\le e_+$ is a closed interval strictly inside $(-\pi,\pi)$. Choose collars between it and the seam where $V-e_+$ has a fixed positive lower bound. Extend $V$ smoothly outside $(-\pi,\pi)$ to a confining potential $\widetilde V$ with no additional well or critical value below $e_+$. Let $\widetilde P_\hbar=-\hbar^2\partial_k^2+\widetilde V$ on the line and let $\mu_j$ be its eigenvalues below $e_+$.

Dirichlet and Neumann realizations on $(-\pi,\pi)$ bracket every twisted realization in quadratic-form order, because
\[
 H_0^1(-\pi,\pi)\subset
 \{g\in H^1:g(\pi)=e^{-2\pi iy}g(-\pi)\}
 \subset H^1(-\pi,\pi).
\]
The bracket is independent of the twist. Cutting the line at the same seam gives the same brackets for $\widetilde P_\hbar$ below $e_+$; the two exterior pieces have a strictly higher form lower bound.

\emph{Exponential independence of the boundary condition.} Choose a real even Lipschitz weight $\phi$, zero on a fixed neighbourhood of the allowed set for energies at most $e_+$, increasing through the forbidden collars, and constant with the same positive value near both endpoints. Take that neighbourhood slightly larger than the allowed set, so that all transitions of $\phi$ occur in a fixed strictly forbidden region. Require $|\phi'|^2\le(V-e_+)/2$ wherever $\phi$ is nonzero. Multiplication by $e^{\phi/\hbar}$ preserves each twisted boundary condition; boundary terms also vanish for the Dirichlet and Neumann realizations. The eigenfunction equation therefore gives, for an eigenfunction $g$ with eigenvalue $E$,
\[
 \int\bigl(\hbar^2|(e^{\phi/\hbar}g)'|^2
 +(V-E-|\phi'|^2)e^{2\phi/\hbar}|g|^2\bigr)=0.
\]
The negative part lies where $\phi=0$ and is bounded by normalization. Positivity on the forbidden collars bounds the weighted mass and scaled derivative there uniformly for $E\le e_+$. On an outer collar where $\phi\ge c>0$, the unweighted mass and scaled derivative are at most $Ce^{-c'/\hbar}$. Use the same argument for the confining line realization. Cutting off before the seam now changes norms, scalar products and Rayleigh quotients by at most $C\hbar^{-M}e^{-c'/\hbar}$. The subspace below $e_+$ has dimension $O(\hbar^{-1})$, so the estimates remain exponential after summing over that subspace. Cut-off line eigenfunctions lie in $H_0^1(-\pi,\pi)$ and bound the Dirichlet levels from above; cut-off Neumann eigenfunctions extend by zero to the line, where $\widetilde V=V$ on their support, and bound the levels $\mu_j$ from above. Dirichlet/Neumann bracketing and min--max therefore give
\begin{equation}\label{eq:twisttunnel}
 |e_j(\kappa,y)-\mu_j(\kappa)|\le Ce^{-c''/\hbar},
\end{equation}
uniformly in $y$, for all indices in the smaller energy window $\mu_j\le3.5$. The same bracketing excludes extra levels there.

\emph{Uniform spacings within the line well.} For $0<e<e_+$, let $k_e=\arccos(1-e/2)$ and use the action
\begin{equation}\label{eq:action}
 \mathcal A(e)=\frac1\pi\int_{-k_e}^{k_e}
                    \sqrt{e-2+2\cos k}\,dk.
\end{equation}
It extends smoothly to zero, with $\mathcal A(e)=e/2+O(e^2)$. Its derivative satisfies $0<\ell_-\le\mathcal A'\le\ell_+$ on $[0,e_+]$ (decrease $e_+$ slightly if necessary). Away from zero the two turning points are simple. The one-dimensional Bohr--Sommerfeld construction with the two Airy connection phases gives
\begin{equation}\label{eq:BS}
 \mathcal A(\mu_j)=\hbar(j+1/2)+O(\hbar^2)
\end{equation}
with an error uniform on compact regular energy intervals. The regular Mathieu action and its higher-order quantum actions are recorded in Hezari--Zelditch~\cite[Section~3.4]{HezariZelditch}; the uniform one-dimensional quantization theorem is Helffer--Robert~\cite[Theorem~3.14 and Corollary~3.15]{HelfferRobert}.

To include the bottom, use the nondegenerate minimum $V(k)=k^2+O(k^4)$. For the fixed confining scalar well, apply Helffer--Robert~\cite[Theorem~4.5 and Remark~4.9]{HelfferRobert} near the Morse minimum, and~\cite[Theorem~3.14 and Corollary~3.15]{HelfferRobert} on overlapping compact regular intervals. The near-minimum expansion is uniform above a fixed finite index. For each of the remaining finitely many indices, oscillator perturbation gives $\mu_j=(2j+1)\hbar+O_j(\hbar^2)$; since $\mathcal A(e)=e/2+O(e^2)$, it gives the same action formula. Sturm ordering fixes the index from the ground state and identifies the indexing on the overlap. The principal action is smooth through zero and the Schr\"odinger subprincipal symbol vanishes. Consequently
\[
 \sup_j|\mathcal A(\mu_j)-\hbar(j+1/2)|\le C\hbar^2
\]
throughout the chosen fixed window below the separatrix.

These results are applied to $\widetilde P_\hbar$. Its potential $\widetilde V$ agrees with $V$ on the circle and remains above $e_+$ outside it. It has one well below $e_+$, a nondegenerate minimum, compact low-energy sublevels and regular closed orbits at positive energies. The normalized action gives period $2\pi$. These are the scalar one-well hypotheses being used, with no $q$ parameter. Uniformity is over $j$ up to a constant$/\hbar$, not just fixed $j$.

Subtracting adjacent action equations gives $\mathcal A(\mu_{j+1})-\mathcal A(\mu_j)=\hbar+O(\hbar^2)$. The bounds $0<\ell_-\le\mathcal A'\le\ell_+$ then give
\[
 \mu_{j+1}-\mu_j\ge\hbar/(2\ell_+)\ge c_M\hbar
\]
for small $\hbar$ and $c_M\le1/(2\ell_+)$. This step is valid even though the individual remainders need not vary smoothly with $j$. Inverting the action also gives $\mu_j\le C_M(j+1)\hbar$ and at least $c_1/\hbar$ levels below $3.5$, after fixing a sufficiently small $c_1$. Finally~\eqref{eq:twisttunnel} is less than $c_M\hbar/8$ for small $\hbar$. This proves~\eqref{eq:Mladder} and the common-window assertion.
\end{proof}

The analytic work imported in this argument is the standard scalar one-well quantization theorem, including its Morse-minimum version; the twist conclusion is the explicit bracketing and forbidden-region argument above. Neither cited source is described as stating the lattice-offset statement verbatim. In particular the input does not assert stability for an arbitrary $q$-dependent quantized symbol.

\subsection{Classical characteristic values}
The same Fourier transformation identifies the union of the offset spectra with the Floquet spectrum of the Mathieu equation. Let $a_D,b_D$ be its characteristic values. Translation by $\pi/2$ changes the sign of the Mathieu parameter if required by convention. Then
\begin{equation}\label{eq:Mathieugap}
 G_D(\kappa)=\frac\kappa4
       \bigl(a_D(4/\kappa)-b_D(4/\kappa)\bigr)>0.
\end{equation}
The endpoints occur at offsets zero and $1/2$. The classical fixed-$D$ asymptotics are
\[
 G_D(\kappa)=\frac{2\kappa^{1-D}}{((D-1)!)^2}
                (1+O_D(\kappa^{-1}))\quad(\kappa\to\infty),
 \qquad G_D(\kappa)\sim2\sqrt\kappa\quad(\kappa\to0).
\]
See the standard Mathieu expansions~\cite{DLMF}. The leading phase-space count below the separatrix is $8/(\pi\sqrt\kappa)$, obtained by integrating the width of $\{\kappa j^2+2-2\cos k\le4\}$. None of these fixed-index formulas replaces the uniform argument above.

\label{end:mathieu}

\section{The conserved polynomial and its first-order structure}\label{app:remainder}
These results explain why a norm bound alone misses the edge ladder. The global proof above selects the branch directly in energy space.

\subsection{Residue regrouping of Paper A}
Paper~A, Proposition~3.3, gives
\begin{equation}\label{eq:polyidentity}
 P(H)=V+S^q+S^{-q}+R_q,\quad
 V(j)=2\Lambda\cos2\pi(q\theta+j\delta),\quad
 \norm{R_q}\le C_hq^{5/2}e^{hq}|\delta|.
\end{equation}
The remainder is self-adjoint of range $q$. Under the unitary residue map
$(\mathcal Uu)_r(m)=u(r+mq)$,
\begin{equation}\label{eq:regroup}
 V+S^q+S^{-q}
 =\mathcal U^*\left(\bigoplus_{r=0}^{q-1}
 H_{\Lambda,q\delta,q\theta+r\delta}\right)\mathcal U.
\end{equation}
Indeed shifts by $q$ preserve a residue class, and the potential there is
$2\Lambda\cos2\pi(q\theta+r\delta+mq\delta)$.
This is an immediate regrouping of the companion result, not a new operator estimate.

All summands have the same phase-independent irrational spectrum $\Sigma'_q$. Spectral mapping gives
\[
 d_{\rm H}(P(\Sig),\Sigma'_q)\le\norm{R_q}.
\]
If $\mu$ is the DOS measure of $H$ and $N'$ the IDS of a summand, then
\[
 N'(t-\norm{R_q})\le\mu(\{E:P(E)\le t\})
                       \le N'(t+\norm{R_q}).
\]
To see this, compress to long boxes, apply min--max, and pay the rank $O(q)$ boundary discrepancy between a compressed polynomial and the polynomial of a compression. Every residue class has density $1/q$, so there is no extra factor $q$ in the DOS.

For $b\ge c>0$ choose $h-L<c$ to make the available error small relative to $\Lambda$. Absolute smallness follows on $b\ge L+\varepsilon$ with $h-L<\varepsilon$. If $B_q=C_hq^{5/2}e^{hq}|\delta|$ denotes the \emph{upper bound}, then
\[
 \frac{B_q}{2\hbar}=\frac{C_h}{4\pi}q^{3/2}e^{(h-L/2)q}\to\infty.
\]
Thus that bound supplies no resolution of the level spacing. This is not a lower bound on the actual remainder and proves no impossibility for a finer analysis. The transformed frequency is exactly $q\delta=q^2(\alpha-p/q)$. Its reciprocal has a continued-fraction carry in general; no semigroup renormalization identity is asserted.

\subsection{An explicit second-order remainder}
Write $P(E)=\sum_nc_nE^n$ and, for the frozen full-line operator $B_t$, define
\begin{equation}\label{eq:Gdefinition}
 \mathcal G_t=\frac12\sum_nc_n\sum_{a+b=n-1}
 \bigl([B_t^a,X]D_tB_t^b+B_t^aD_t[X,B_t^b]\bigr).
\end{equation}
It is periodic, real symmetric, and has range at most $q-1$. Let $\Op^{\rm row}(\mathcal G)$ freeze its entries at the row phase $\phi_j=\theta+j\eps$.

\begin{proposition}[First-order expansion]\label{prop:remainder}
For fixed $h>L$, eventually and uniformly in $\theta$,
\begin{equation}\label{eq:R2}
 R_q=\eps\Op^{\rm row}(\mathcal G)+R_2,
 \qquad \norm{R_2}\le C_hq^4e^{hq}\delta^2.
\end{equation}
Consequently on $b\ge L/2+\xi$, with an appropriate fixed $h-L<\xi$,
$\norm{R_2}=o(\hbar)$ uniformly.
\end{proposition}
\begin{proof}
Differentiating the rational magic formula $P(B_t)=S^q+S^{-q}+c(t)I$ gives $dP(B_t)[D_t]=c'(t)I$. Symmetrize the insertion of $X-j$ into the polynomial derivative. The commutators left over are~\eqref{eq:Gdefinition}, and
\[
 dP(B_t)[(X-j)D_t]=c'(t)(X-j)+\mathcal G_t.
\]
The first term has zero row $j$.

As in the row proof of Paper~A, replace the rational potential only on $[j-q,j+q]$ by the actual one. The resulting diagonal perturbation is $W=\eps(X-j)D_t+W_2$ on this window, with
$\norm W\le4\pi\lambda|\delta|$ and $\norm{W_2}\le4\pi^2\lambda\delta^2$. The first derivative on row $j$ sees no entries outside the window.

Choose the fixed smooth energy cutoff $f_q=\chi P$ from the proof of Paper~A, Proposition~3.3. Its support lies in the derivative neighbourhood, and it equals $P$ near both comparison spectra. Rational spectral continuity and $\norm W\to0$ justify that assertion uniformly in the row. The derivative bounds imply $\norm{f_q}_{H^r}\le C_{r,h}q^re^{hq}$. Fourier inversion and Duhamel give
\[
 \norm{df_q(B)[W_2]}\le Cq^2e^{hq}\norm{W_2},\quad
 \norm{f_q(B+W)-f_q(B)-df_q(B)[W]}
       \le Cq^3e^{hq}\norm W^2.
\]
Thus each row of the error in~\eqref{eq:R2} has $\ell^2$ norm at most $Cq^3e^{hq}\delta^2$. Each entry has that bound, and both a row and a column have at most $2q+1$ entries. Direct banded Schur summation gives the claimed $q^4$ bound; no symmetry of row freezing is assumed.

Finally
\[
 \frac{\norm{R_2}}{2\hbar}\le C_hq^3e^{(h-L/2-b)q},
\]
which decays for the stated fixed margin. This proves sufficiency, not necessity at or below $b=L/2$.
\end{proof}

Although every frozen $\mathcal G_t$ is symmetric, its row quantization need not be. Taking the Hermitian part of the full first-order comparison costs no more than $\norm{R_2}$, since $P(H)$ is self-adjoint. Such a comparison still concerns $P(H)$ and does not recover individual inverse branches or their IDS labels by itself.

\subsection{Orbital and frame entries}
For simple fibre eigenvalues $E_x$ and normalized eigenvectors $u_x$, write $D=\partial_tB$, $K=\partial_kB$. Divided differences with repeated nodes are understood by continuity.

\begin{proposition}[Band-frame symbol]\label{prop:symbol}
In the band frame,
\begin{equation}\label{eq:divdiff}
 \mathcal G_{xy}=\frac{iq}{2}\sum_zP[E_x,E_z,E_y]
                    (D_{xz}K_{zy}-K_{xz}D_{zy}).
\end{equation}
In particular
\begin{align}
 \mathcal G_{xx}
 &=-qP'(E_x)\sum_{z\ne x}
        \frac{\operatorname{Im}(D_{xz}K_{zx})}{E_x-E_z},
 \label{eq:orbital}\\
 \mathcal G_{xy}
 &=-iq\bigl(c'(t)\inner{u_x}{\partial_ku_y}
       +2\sin k\inner{u_x}{\partial_tu_y}\bigr),\quad x\ne y.
 \label{eq:frameentries}
\end{align}
Every simple phase-union edge has vanishing off-diagonal entries in~\eqref{eq:frameentries}.
\end{proposition}
\begin{proof}
The fibre map obeys $[A,X](k)=-iq\partial_kA(k)$. Differentiate the powers in~\eqref{eq:Gdefinition} and insert the spectral decomposition between their factors. The resulting scalar sums are precisely the second divided differences, proving~\eqref{eq:divdiff}. All fibre eigenvalues have the same $P$-value $2\cos k+c(t)$. Hence a divided difference at three distinct nodes is zero, while
$P[E_x,E_z,E_x]=P'(E_x)/(E_x-E_z)$.
The diagonal commutator gives~\eqref{eq:orbital}. For $x\ne y$ only the repeated nodes $z=x,y$ remain. Hellmann--Feynman gives
$P'(E_x)D_{xx}=c'(t)$ and $P'(E_x)K_{xx}=-2\sin k$; differentiation of eigenvectors yields~\eqref{eq:frameentries}. At a phase-union edge, $\sin k=0$ and $c'(t)=0$. Some other critical points are saddles; the converse characterization is not asserted.
\end{proof}

For a simple edge $E_a$ define
\[
 S_a=q\sum_{z\ne a}\frac{\operatorname{Im}(D_{az}K_{za})}{E_a-E_z},
 \qquad \sigma_a=\eps\mathcal G_{aa}=-\eps P'(E_a)S_a.
\]
The E-space orbital correction is $-\eps S_a$, a recovery of the Rammal--Wilkinson term~\cite{RammalBellissard,Wilkinson1987,PST}. It is only the orbital part of the edge displacement: the Landau zero-point term is proportional to $|\eps|$ and supplies a separate one-sided contribution. The safe bound is $|S_a|\le q\norm D\norm K/g_a=O(g_a^{-1})$, where $g_a$ is the simple-band isolation. At the top edge this is subexponential. No uniform $O(1)$ bound for all nearly touching bands, no asymptotic $P'\sim qe^{hq}$, and no nonzero universal orbital lower coefficient are claimed.

At an edge with $P(E_a)=\sigma R$, the signed shifted P-space centre is $\sigma R+\sigma_a$, not $\sigma(R+\sigma_a)$. A single-band spectral correspondence at a negative edge would require that convention and its own branch proof. The main theorem uses only $\sigma=+1$.

\subsection{Closed gaps: the recovered Dirac model}
At a central cone with $q\equiv0\pmod4$, $P(0)=R$ and $P'(0)=0$. Let $\Pi_0$ be the rank-two crossing projection. The degenerate polynomial calculus gives
\begin{equation}\label{eq:Diracblock}
 \Pi_0\mathcal G\Pi_0
   =\frac{iqP''(0)}4[\Pi_0D\Pi_0,\Pi_0K\Pi_0].
\end{equation}
Chambers' quadratic expansion has perpendicular Pauli velocities
\[
 v_k=\sqrt{2/|P''(0)|},\qquad
 v_t=\sqrt{8\pi^2q^2\Lambda/|P''(0)|}.
\]
Thus~\eqref{eq:Diracblock} has eigenvalues $\pm2\pi q^2\sqrt\Lambda$. After multiplication by $\eps$, they are $\pm\sgn(\delta)\hbar$. The block is nonzero; cancelling an inverse gap in a simple-band formula does not make its degenerate extension vanish.

Quantizing the two linear Pauli directions and squaring gives oscillator levels $(2j+1)\hbar$ with the spin correction $\pm\hbar$. Consequently the leading Dirac model has one zero mode and pairs
\[
 E_D=\pm2\sqrt{\hbar D/|P''(0)|},\qquad D\ge1.
\]
This is the known Dirac Landau fan of Harper systems; see Claro~\cite[(4)--(6)]{Claro} and Bellissard--Barelli~\cite[Section~1.3]{BellissardBarelli}. Since $P$ is even, the paired model energies have the same P-value. Exact degeneracy of a finite scalar Jacobi matrix, an exact zero eigenstate of the full irrational operator, and an all-order anharmonic formula do not follow. Its satellite-gap transfer remains a separate problem. None of this degenerate model is needed for the top simple edge proof.

\label{end:remainder}

\section{Two classical limits with their quantifiers}\label{app:limits}
\subsection{A fixed-curvature Mathieu limit}
\begin{proposition}[Fixed $\kappa$ limit and the upper-edge sign]\label{prop:fixedkappa}
Fix $\kappa>0$. Let $\Lambda\to\infty$ and
$\omega=\sqrt\kappa/(2\pi\sqrt\Lambda)>0$ be irrational. For any phase put
\[
 T_\Lambda=2\Lambda+2-H_{\Lambda,\omega,\theta}.
\]
On a bounded energy window whose upper endpoint is in a model gap, its spectrum converges in Hausdorff distance to
$\bigcup_y\sigma(\hM_{\kappa,y})$, with error $O_\kappa(\Lambda^{-1/6})$. Every fixed model gap inside the window is approached by a gap of the same limiting width. Its upper-edge label for $H_{\Lambda,\omega,\theta}$ is $-D$, with IDS $1-D\omega$. The reflected bottom gap has label $+D$.
\end{proposition}
\begin{proof}
The potential of $T_\Lambda$ is $W(m)=2\Lambda(1-\cos2\pi(\theta+m\omega))$, and its kinetic term is $2-S-S^*$. Its bottoms are $m_k=(k-\theta)/\omega$, separated by $1/\omega$. Their offsets modulo integers are dense, because $1/\omega$ is irrational. Near a bottom,
\[
 |W(m)-\kappa(m-m_k)^2|
 \le\frac{\kappa^2(m-m_k)^4}{12\Lambda}.
\]
At distance at least $M$ from every bottom, $W\ge(4/\pi^2)\kappa M^2$.

Take $M=\Lambda^{1/6}$. A low-energy normalized model eigenvector can be cut off within distance $M$ with residual $O(M^{-1})$, uniformly over its compact offset circle and the finitely many levels in the fixed window. Choose a true well with offset error $\eta\le M^{-2}$, after integer translation of the offsets. Its phase mismatch contributes $O_\kappa(\eta M)$, and the quartic contributes $O_\kappa(M^4/\Lambda)$. The normalized cutoff is therefore a quasimode of $T_\Lambda$ with error $O_\kappa(\Lambda^{-1/6})$.

For the reverse inclusion choose a squared partition with bounded overlap, with each well cutoff equal to one within radius $M$ and supported within radius $2M$; require every well cutoff and the exterior cutoff to have discrete gradient at most $C/M$. Enlarge supports by one site when evaluating the nearest-neighbour commutators. The operator sum of their squared commutators is $O(M^{-2})$; diagonal terms commute with the cutoffs. The Taylor estimate is used on radius $2M$, which changes only its constant, and the exterior starts at distance at least $M$ from every bottom. A quasimode of $T_\Lambda$ in the fixed window has exterior mass tending to zero, since the exterior form lower bound is $\kappa M^2/\pi^2$. Averaging the squared localized residuals gives at least one well whose normalized residual has size $O(M^{-1})$. Comparing its potential with the quadratic model adds $O_\kappa(M^4/\Lambda)$. This proves the reverse spectral approximation and the asserted rate. Compactness of the window and uniform model tails make both inclusions uniform. Neighbouring model-band witnesses locate the endpoints of each fixed gap.

To label the gap, take $t$ strictly inside a fixed model gap. On the exterior, $T_\Lambda-t$ is positive. Its elimination contributes a Schur correction of norm at most $C/(\kappa M^2)$ on the union of the well windows, including all inter-window terms. The model's own positive exterior tails justify that its large Dirichlet restrictions have exactly $D$ levels below $t$, uniformly in the offset. The quartic and Schur errors are smaller than the fixed distance from $t$ to the union of model spectra. Thus a finite collection of complete windows has exactly $D$ negative directions per well after the errors, by homotopy and min--max. Partially cut boundary windows cost a bounded number for fixed $\Lambda$ and vanish in the box limit. There are $\omega$ wells per site, so $N_{T_\Lambda}(t)=D\omega$. Energy reversal gives
\[
 N_{H_{\Lambda,\omega}}(2\Lambda+2-t)=1-D\omega.
\]
Gap labelling identifies the upper label as $-D$. Reflection gives $+D$ at the bottom.
\end{proof}

This is a fixed-$\kappa$, fixed-window statement. Its constants need not be uniform as $\kappa\to0$ or $\infty$, or as $D$ grows. It therefore does not replace Appendix~\ref{app:comparison}. It is a self-contained application of classical separated-well ideas; no firstness is claimed for the scaling mechanism.

\subsection{Fixed-period strong coupling and the Sudler product}
Write $S_N(\omega)=\prod_{j=1}^N|2\sin(\pi j\omega)|$, with $S_0=1$. Rational labels in the next proposition mean minimal representatives modulo $Q$.

\begin{inputbox}
\begin{classicalproposition}[Classical fixed-period perturbation calculation]
Fix reduced $p/Q$ with odd $Q$, and $1\le d=|m|\le(Q-1)/2$. The width of its rational phase-union gap obeys
\begin{equation}\label{eq:Sudler}
 w_m(\lambda)=2\lambda^{1-d}S_{d-1}(p/Q)^{-2}
                         (1+O_{p,Q,m}(\lambda^{-1}))
 \quad(\lambda\to\infty).
\end{equation}
\end{classicalproposition}
\end{inputbox}
\begin{proof}
Divide the $Q$-site fibre by $\lambda$ and put $\epsilon=1/\lambda$. The unperturbed diagonal is $f(\phi_j)=2\cos2\pi\phi_j$, $\phi_j=jp/Q+\theta$, with hopping $\epsilon$. A crossing between sites at shortest lattice distance $d$ has $2\phi_j=-dp/Q\pmod1$. It contains exactly two resonant levels with opposite nonzero slopes; all other levels are separated by a positive constant for this fixed rational.

Eliminate the other sites in a small fixed crossing neighbourhood. The off-diagonal Feshbach entry is
\[
 T=\epsilon^d\prod_{i=1}^{d-1}
     (E-f(\phi_{j+i}))^{-1}(1+O(\epsilon^2))
       +O(\epsilon^{Q-d}).
\]
The first path is the unique shortest arc; the second arc has relative size $\epsilon^{Q-2d}=O(\epsilon)$. Backtracking adds at least two hops. The estimates and their local energy and phase derivatives follow from the convergent nonresonant resolvent series.

At the corrected crossing $\theta_0$, the derivative of the sum of the two diagonal entries is $O(\epsilon^2)$, while their difference has nonzero derivative. This small common slope assertion is used at $\theta_0$, not throughout the fixed neighbourhood. Retain the common quadratic Taylor term there. Optimizing the two avoided-crossing branches in $\theta$ then gives width $2|T|(1+O(\epsilon))$; the phase displacement is $O(|T|)$ and its quadratic contribution is a relative $O(\epsilon)$ or smaller error. Dependence of the effective matrix on energy has derivative $O(\epsilon^2)$ and does not change that leading term.

At the crossing,
\[
 |E_*-f(\phi_{j+i})|
 =4|\sin(\pi ip/Q)\sin(\pi(d-i)p/Q)|,
\]
so their product is exactly $S_{d-1}(p/Q)^2$. Worst Bloch phase changes the answer only by the complementary-arc relative error. Finite classical ordering excludes any unrelated gap closing outside the finitely many equivalent crossing neighbourhoods; their gap index satisfies $k\equiv\pm dp\pmod Q$. This fixes $|m|=d$. Multiplying the normalized width by $\lambda$ proves~\eqref{eq:Sudler}.
\end{proof}

The sine-product structure is classical. Fedotov~\cite[Theorem~2.5]{Fedotov} gives a related small-coupling lower bound with a finite label range and its own frequency-dependent smallness hypotheses. Its printed coefficient is half the dual fixed-rational asymptotic coefficient here. This calculation supplies no uniform remainder as $Q$ grows and no irrational label anchoring by itself. Applying it recursively to the transformed operator is the heuristic Law R, not a theorem proved in this paper.

\label{end:limits}

\section{Results from the companion papers}\label{app:companions}
This is the complete companion-import list for the present manuscript. The versioned references are Paper~A~\cite{PaperA}, v2.0.1; Paper~B~\cite{PaperB}, v1.0.0; and Paper~C~\cite{PaperC}, v1.0.0. Numbers and source labels below refer to these versions, whose DOIs are in the bibliography.

\begin{enumerate}
\item \textbf{A, Lemmas 2.5--2.6, \texttt{prelim.tex}, labels \texttt{lem:window}, \texttt{lem:ratspec}}. A supported Weyl vector has residual at most $8/M$; rational phase unions approach the irrational spectrum with a square-root frequency bound.

\emph{Use here:} Equation~\eqref{eq:continuity}; the surplus conversion; the derivative domain; Corollary~\ref{cor:derivatives}.

\item \textbf{A, Proposition 3.3(1), \texttt{polynomial.tex}, label \texttt{prop:poly}}. $P(H)=V+S^q+S^{-q}+R_q$, with range-$q$ self-adjoint remainder bounded by $C_hq^{5/2}e^{hq}|\delta|$.

\emph{Use here:} Equations~\eqref{eq:polyidentity}--\eqref{eq:regroup}; interpretation of the coarse norm error.

\item \textbf{A, Proposition 3.3(2) and proof Steps 1--3}. Near-spectrum derivative bounds $|P^{(r)}|\le C_{r,h}q^re^{hq}$; locally modified rational comparison and a fixed smooth cutoff.

\emph{Use here:} Top-band length and width bounds; inverse curvature; $m_0$; Proposition~\ref{prop:even}; Proposition~\ref{prop:remainder}.

\item \textbf{A, Theorem A}. Typical moment and fastest upper Abel exponents $\max(0,1-L/\beta)$; all-phase statement for $\beta=\infty$. A gives the finite-positive-$\beta$ law used here; its $\beta=\infty$ transport and packing endpoint is due to Jitomirskaya--Zhang~\cite[Corollary~3]{JZ2022}, and its $\beta=0$ fastest endpoint to Damanik--Tcheremchantsev~\cite{DT2007}.

\emph{Use here:} Section~\ref{sec:diagram}, including the typical-transport entry and the all-phase infinite-$\beta$ endpoint.

\item \textbf{A, Theorem C}. Both packing dimensions of every nonzero spectral measure have the same value under its a.e.-phase hypotheses, and at every phase for $\beta=\infty$ (the Jitomirskaya--Zhang endpoint~\cite[Corollary~3]{JZ2022}).

\emph{Use here:} The spectral-measure packing statement in Section~\ref{sec:diagram}.

\item \textbf{B, Lemmas 2.4--2.5, \texttt{prelim.tex}, labels \texttt{lem:growth}, \texttt{lem:ratprod}}. Uniform real-energy transfer growth and its rational-product expansion.

\emph{Use here:} Appendix~\ref{app:complex} proves the needed complex extension. The extension is proved here, rather than quoted as the original statement.

\item \textbf{B, Theorem A}. The typical fastest-exponent law for supercritical type I operators.

\emph{Use here:} Background attribution in Section~\ref{sec:diagram}; the AMO table itself uses A.

\item \textbf{C, Corollary 7.2, \texttt{sharp.tex}, label \texttt{cor:closedgap}}. For even $q$, $P$ is even and $P(0)=(-1)^{q/2}(2\lambda^q+2)$.

\emph{Use here:} Proposition~\ref{prop:even}; the central-cone discussion in Appendix~\ref{app:remainder}.

\item \textbf{C, Theorem 7.5, \texttt{sharp.tex}, label \texttt{thm:quarter}}. At even $q_n$, a polynomially small fastest tail at radius $c_\eta Te^{-hq_n/2}$ for $b_n\ge h/2+(h-L)$ and $e^{(h/2+(h-L))q_n}\le T\le|\delta_n|^{-1}$; nearby turning-phase clause.

\emph{Use here:} The parity-qualified phase-$1/4$ lower bound and transport geometry in Section~\ref{sec:diagram}.

\item \textbf{C, Corollary 5.6, \texttt{upperlaw.tex}, label \texttt{cor:typeI}}. For AMO, $\alpha_u^{\rm unif}\le\beta/(\beta+L)$ below $\beta=L$ and $\le1-L/(2\beta)$ above, including $\beta=L$.

\emph{Use here:} Every upper transport entry in Table~\ref{tab:phase}; the endpoint $\beta=0$.

\item \textbf{C, Corollary 7.6, \texttt{sharp.tex}, label \texttt{cor:exactunif}}. If $\beta_e=\beta\in[L,\infty)$, the uniform, worst-phase and phase-$1/4$ exponents equal $1-L/(2\beta)$.

\emph{Use here:} Exact entries at $x\ge2$ with the stated parity hypothesis.

\item \textbf{C, Proposition 6.4, \texttt{amoproof.tex}, label \texttt{prop:Camo}}. Uniform confinement radius bounded by a constant times $\sqrt{\sigma_n}/|\delta_n|$ plus its specified collars.

\emph{Use here:} The geometric comparison in Section~\ref{sec:diagram}; not a new confinement estimate.
\end{enumerate}

In A v2.0.1, Lemma~2.5 is the window lemma and Lemma~2.6 the rational-spectrum lemma; in C v1.0.0, Corollary~7.2 is the closure corollary and 7.3 its priority remark. The source labels guard against numbering changes between versions.

For clarity, the external analytic inputs are separate from this table. Chambers' formula is due to Chambers~\cite{Chambers1965}; the periodic polynomial identity also follows from the standard magic formula~\cite[Theorem~3.1]{DKS2010}. On the spectrum, the Lyapunov exponent equals $\log\lambda$, by Bourgain--Jitomirskaya~\cite[Corollary~2]{BourgainJitomirskaya2002}, after matching the coupling convention. The IDS facts are standard AMO background, as used in A/B and reviewed in~\cite{AYZdry}. The rational core measure is Avron--van Mouche--Simon~\cite{AvMS}; the all-irrational noncritical measure is Jitomirskaya--Krasovsky~\cite{JKmeasure}. The one-well quantization input is explicitly scoped in Appendix~\ref{app:mathieu}. The Aubry--Perron fact needed at the reference fibre is proved in Lemma~\ref{lem:Wannier}. No unproved general Harper perturbation lemma, all-band reduction, or dynamical localization assumption is imported.

For the new gap theorem the companion inputs are A, Proposition~3.3(2), and B, Lemmas~2.4--2.5, as extended in Appendices~\ref{app:complex} and~\ref{app:below-family}. The one-band range extension is proved in Lemma~\ref{lem:belowrange}; the new inverse, locality, self-energy and model estimates are proved in Appendices~\ref{app:below-resolvent}--\ref{app:below-model}. The additional lower bound uses only Liu--Shi, Theorem~1.3, and the duality proof in Section~\ref{sec:below-literature}. Paper~C supplies the separately labelled transport picture and even-denominator surplus variant. No transport theorem is imported into the new fixed-$D$ proof.

\label{end:companions}

\bibliographystyle{plainurl}
\bibliography{references}
\end{document}
